\documentclass[11pt,a4paper]{article}

\usepackage{amsmath,amssymb,amsthm}
\usepackage{geometry}
\usepackage{algorithm}
\usepackage{algpseudocode}
\usepackage{enumitem}
\usepackage{booktabs}
\usepackage{hyperref}

\geometry{a4paper,margin=25mm}

\theoremstyle{plain}
\newtheorem{theorem}{Theorem}[section]
\newtheorem{proposition}[theorem]{Proposition}
\newtheorem{lemma}[theorem]{Lemma}
\newtheorem{corollary}[theorem]{Corollary}
\theoremstyle{definition}
\newtheorem{definition}[theorem]{Definition}
\theoremstyle{remark}
\newtheorem{remark}[theorem]{Remark}

\newcommand{\R}{\mathbb{R}}

\title{Simultaneous Classification and Solution of Linear Optimization Problems by the Dual Simplex Method: The Slope--Intercept Two-Phase Dual Simplex Method}
\author{Kan Enomoto\thanks{Website: \url{https://enomoto-math-labo.com/}}\\
Independent Researcher, Usa, Oita, Japan\\
\texttt{enomotokan@gmail.com}\\
ORCID: \href{https://orcid.org/0009-0001-9485-0393}{0009-0001-9485-0393}}
\date{}

\begin{document}
\maketitle

\begin{abstract}
The all-slack basis becomes dual feasible without any computation if each nonbasic variable is placed at the bound indicated by the sign of its cost, but only when those bounds are finite. The artificial bounding method removes this obstacle by replacing infinite bounds with $\mp M$, but a fixed numerical $M$ may be too small and causes numerical difficulties. We instead treat all quantities arising in the iterations as affine functions of $M$ and analyze the limit $M\to+\infty$ symbolically. The iterations then split exactly into two phases: Phase~A solves the \emph{slope problem}, which coincides with the auxiliary problem known as Phase~1 of the dual simplex method, and Phase~B solves the \emph{intercept problem} obtained by fixing the slope. This decomposition connects the artificial bounding method with auxiliary-problem Phase~1. On this basis we derive the \emph{slope--intercept two-phase dual simplex method}. Using only dual simplex iterations and a light cleanup, it simultaneously classifies a linear optimization problem as having a finite optimum, being unbounded, or being infeasible, and computes an optimal basic solution when one exists, without a case distinction on the outcome of Phase~A and without re-solving. Standard techniques such as the dual steepest-edge rule and the bound flipping ratio test apply unchanged. Experiments on 207 instances show that our implementation is competitive with HiGHS, CLP and SoPlex and distinguishes infeasibility from unboundedness at little extra cost.
\end{abstract}

\noindent\textbf{Keywords:} linear optimization $\cdot$ dual simplex method $\cdot$ infeasibility detection $\cdot$ unboundedness $\cdot$ big-$M$ method $\cdot$ artificial bounds

\medskip
\noindent\textbf{Mathematics Subject Classification (2020):} 90C05 $\cdot$ 65K05

\section{Introduction}
\label{sec:intro}

The simplex method for linear optimization (LP; also called linear programming) is a classical solution method that has been studied for many years in both theory and implementation. To actually run the simplex method, one must start from an initial basic solution that is \emph{feasible}; that is, one must choose an invertible $m\times m$ submatrix (basis) $A_B$ of the constraint matrix $A$ such that, when the nonbasic variables are fixed at their upper or lower bounds, the values of the basic variables
\[
  x_B = A_B^{-1}\bigl(b - A_N x_N\bigr)
\]
all lie within their own bounds.

Constructing such an initial basis is not trivial in general. For an \emph{arbitrary} linear optimization problem consisting of equality constraints $Ax=b$ and general range constraints $l\le x\le u$ ($l_j,u_j\in\R\cup\{\pm\infty\}$), no general method is known that provides a feasible initial basis at no computational cost, and auxiliary procedures such as the two-phase method and the big-$M$ method \cite{dantzig1963} have been used (Section~\ref{sec:difficulty}).

On the other hand, if one focuses on dual feasibility, then for the all-slack basis consisting only of slack variables with cost $0$, the reduced costs coincide with $\bar c_j=c_j$; hence, merely by placing the nonbasic original variables as
\[
  c_j\ge 0 \ \Longrightarrow\ x_j=l_j,\qquad
  c_j< 0 \ \Longrightarrow\ x_j=u_j,
\]
one obtains a dual feasible initial basis \emph{without any computation}. This placement, however, requires the bound at which each variable is placed to be finite. In practical problems many variables have bounds that are finite on both sides, so the variables violating this requirement are often limited to a few that have an infinite side \cite{koberstein2005}. Moreover, computational experiments on real-world problems have reported that the dual simplex method was on average about twice as fast as the primal simplex method \cite{bixby2002}. For this reason, it is of practical significance to be able to carry out, within the framework of the dual simplex method, the complete classification into ``infeasible, unbounded, finite optimum'' that the two-phase primal simplex method provides. In fact, major solvers sometimes return an indistinguishable verdict ``infeasible or unbounded,'' and re-solving with different settings is used to distinguish the two. For example, Gurobi recommends, for this verdict, re-optimizing with the dual reductions of presolve disabled \cite{gurobi}. Furthermore, Gurobi's verdict ``unbounded'' only indicates the existence of a direction along which the objective function can be improved indefinitely, says nothing about the existence of a feasible solution, and Gurobi asks users to re-optimize with the objective function set to $0$ in order to verify feasibility \cite{gurobi}. CPLEX has a similar verdict \cite{cplex}.

The artificial bounding method, which removes this obstacle by replacing infinite bounds with a large constant $M$, is known as one variant of Phase~1 of the dual simplex method and corresponds to the big-$M$ method applied to the dual problem \cite{wolfe1965,koberstein2005,koberstein2007,huang2021}. However, as long as $M$ is fixed as a number, the same problems as in the primal big-$M$ method remain: there is no guarantee that $M$ is large enough, and the introduction of $M$ causes numerical instability. Indeed, Koberstein \cite{koberstein2005} describes an implementation that, if nonbasic variables remain at artificial bounds upon termination, multiplies $M$ by a constant and restarts, and declares unboundedness once a threshold is exceeded; he then excludes this method from his final implementation on the grounds that an appropriate $M$ depends strongly on the problem, that too small an $M$ leads to repeated restarts, and that too large an $M$ causes numerical problems.

As methods that can distinguish the three cases without requiring Phase~1, the criss-cross method \cite{zionts1969,terlaky1985} and the self-dual parametric simplex method \cite{dantzig1963,vanderbei2020} are known. Both can start from an arbitrary basis, but they combine primal and dual pivots and are not the dual simplex method itself. Consequently, implementation techniques developed for the dual simplex method, such as the steepest-edge rule and BFRT, cannot be used as they are. As for the criss-cross method, it has been stated that, in contrast to the many efficient implementations of the simplex method, no practically efficient criss-cross method is known, and its existence has been listed as an open problem \cite[p.~389]{fukuda1997}. In addition, the default LP algorithms of major solvers are the dual simplex method (HiGHS, CPLEX) or a method that runs the dual simplex, interior-point and primal simplex methods concurrently (Gurobi) \cite{huangfu2018,cplex,gurobi}, and these methods are not used.

In this paper we analyze the dual simplex method applied to the bounded problem symbolically in the limit $M\to+\infty$, without fixing $M$ as a number. We show that the iterations then split into two phases that solve the slope problem and the intercept problem in turn, and on this basis we derive a two-phase dual simplex method (the slope--intercept two-phase dual simplex method) in which the steepest-edge rule, the Devex rule and the Bound Flipping Ratio Test (BFRT) can be used as they are.

This paper is organized as follows. Section~\ref{sec:contribution} distinguishes known results from the contributions of this paper; Section~\ref{sec:prelim} reviews the notation of the dual simplex method with bound constraints, and Section~\ref{sec:difficulty} reviews known methods for constructing an initial basis. In Section~\ref{sec:method} we bound the unbounded variables by $M$ and give a trivial dual feasible initial basis. In Section~\ref{sec:symbolic} we analyze the behavior as $M\to+\infty$ symbolically and establish the classification of the original problem into the three cases, the cleanup, the non-return to infinite bounds, and anti-cycling. In Section~\ref{sec:two-phase} we show that the iterations split into two phases that solve the slope problem and the intercept problem and compare this with existing two-phase dual simplex methods, and in Section~\ref{sec:dual-view} we interpret this from the side of the dual problem. Section~\ref{sec:summary} summarizes the slope--intercept two-phase dual simplex method, and Section~\ref{sec:experiments} reports the results of numerical experiments.

\section{Known Results and Contributions of This Paper}
\label{sec:contribution}

Since the content of this paper mixes a reorganization of known results with new theory, we make the distinction between the two explicit here.

\paragraph{Known results.}
The following items are known or follow immediately from known frameworks. In this paper we use them as the foundation of our discussion and, where needed, give proofs in our setting, but we do not regard them as main contributions.
\begin{itemize}
\item The dual simplex method with bound constraints including status $Z$, in which a free variable is kept nonbasic at value $0$ \cite{lemke1954,maros2003} (Section~\ref{sec:prelim}), and the two-phase method, the big-$M$ method, crash procedures and Phase~1 of the dual simplex method for constructing an initial basis (Section~\ref{sec:difficulty}).
\item The artificial bounding method, which replaces infinite bounds with $\pm M$ (Section~\ref{sec:truncation}). This is the dual version of the composite simplex method \cite{wolfe1965}, that is, the big-$M$ method applied to the dual problem \cite{koberstein2005,koberstein2007,huang2021}. Restoring the original bounds of a variable given artificial bounds once it enters the basis, flipping variables with finite bounds on both sides, and combining the method with BFRT have also already been described for implementations using a numerical $M$ \cite{koberstein2005}.
\item The idea of treating $M$ as a non-Archimedean infinity and executing the big-$M$ method by lexicographic comparison of coefficients \cite{sherali1983,cococcioni2021}, and the fact that the big-$M$ method coincides with the two-phase method as $M\to+\infty$ \cite{chvatal1983,bertsimas1997}. Lemma~\ref{lem:affine}, Propositions~\ref{prop:consistency} and \ref{prop:trichotomy}, and the interpretation in Section~\ref{sec:dual-view} correspond to transferring these to the side of the dual problem.
\item The primal push of crossover, which moves superbasic variables with reduced cost $0$ to their bounds \cite{megiddo1991,bixby1994,andersen1996}. The cleanup in Lemma~\ref{lem:cleanup} has the same structure (Remark~\ref{rem:crossover}).
\item Anti-cycling by a lexicographic rule for the ratio test \cite{dantzig1955} (Section~\ref{sec:anticycling}), and the characterization of boundedness by the sign of the objective function on the recession cone \cite{bertsimas1997} (Corollary~\ref{cor:phaseA}).
\item The slope problem itself. It coincides with the auxiliary problem of minimizing the sum of dual infeasibilities, known as Phase~1 of the dual simplex method \cite{fourer1994,kostina2002,maros2003pl,koberstein2005,koberstein2007} (Section~\ref{sec:slope}).
\end{itemize}

\paragraph{Contributions of this paper.}
The contribution of this paper is to treat $M$ in the artificial bounding method symbolically rather than fixing it as a number, to clarify the structure of its iterations, and then to put it into shape as a practical dual simplex method.
\begin{enumerate}[label=(\arabic*)]
\item \textbf{Simultaneous classification and solution by the dual simplex method} (Propositions~\ref{prop:trichotomy} and \ref{prop:alg}): starting from the all-slack basis obtained without computation, and using only iterations of the dual simplex method, we determine exactly whether the original problem has a finite optimal value, is feasible and unbounded, or is infeasible, and, when it has a finite optimal value, we construct an optimal basic solution. Classification and solution are not separate procedures: they are carried out simultaneously within the same two phases, with no case distinction on the outcome of Phase~A and no re-solving of a different problem. The artificial bounding method with a numerical $M$ cannot distinguish an insufficient $M$ from unboundedness and relies on heuristic decisions, while existing two-phase dual simplex methods based on the auxiliary problem stop as soon as dual infeasibility is detected and do not distinguish infeasibility from unboundedness \cite{koberstein2005} (Section~\ref{sec:known-two-phase}).
\item \textbf{Decomposition into slope and intercept} (Propositions~\ref{prop:two-phase} and \ref{prop:phaseB}): the two components of the lexicographic comparison never act simultaneously in the main iterations, and the iterations split, at the level of individual iterations, into Phase~A, which is the dual simplex method applied to the slope problem, and Phase~B, which is the classical dual simplex method applied to the intercept problem obtained by fixing the slope. Because the slope problem appears as the coefficient of $M$ in the artificial bounding method, the artificial bounding method (the big-$M$ method applied to the dual problem) and Phase~1 based on the auxiliary problem (the two-phase method applied to the dual problem), which have so far been treated as separate techniques, are connected (Section~\ref{sec:dual-two-phase}). To the best of our knowledge, no literature has made this correspondence explicit.
\item \textbf{The slope--intercept two-phase dual simplex method} (Sections~\ref{sec:known-two-phase} and \ref{sec:summary}): we give a two-phase algorithm that, after solving in Phase~A the same problem as the existing Phase~1, solves the intercept problem rather than the original problem. By solving the intercept problem directly from the basis of Phase~A regardless of the sign of $z^1$, the same procedure solves the problem and detects infeasibility when $z^1=0$, and distinguishes infeasibility from unboundedness when $z^1<0$. In the case $z^1=0$ where variables remain at artificial bounds, we give, apart from the existing method of re-placing them at their finite sides, a route that obtains an optimal basic solution by the intercept problem and the cleanup while retaining $x^1$, and we understand both in a unified way as differing only in how variables with reduced cost $0$ are placed. The main iterations do not use the lexicographic comparison, and each phase is the dual simplex method applied to an ordinary LP; hence correctness is shown directly (Proposition~\ref{prop:alg}), and the steepest-edge rule, the Devex rule and BFRT can be applied as they are.
\item \textbf{Symbolic cleanup} (Lemma~\ref{lem:cleanup}): we execute the primal push lexicographically in a setting where the step lengths and the basic solution are affine functions of $M$, and construct an optimal basic solution of the original problem without ever fixing $M$ as a number.
\end{enumerate}
On the other hand, the selection rules for the leaving row and the entering variable are essentially unchanged by the separation (Remark~\ref{rem:pricing}), and a comparison of efficiency with existing methods is addressed through the numerical experiments in Section~\ref{sec:experiments}.

\section{Preliminaries: The Dual Simplex Method with Bound Constraints}
\label{sec:prelim}

In this paper we consider linear optimization problems of the following form:
\begin{align}
  \min\quad & c^\top x \label{eq:lp}\\
  \text{s.t.}\quad & Ax = b, \notag\\
  & l_j \le x_j \le u_j, \quad j=1,\dots,n, \notag
\end{align}
where $A\in\R^{m\times n}$, $b\in\R^m$, $c\in\R^n$, $\operatorname{rank}(A)=m$, and $l_j \in \R\cup\{-\infty\}$, $u_j \in \R\cup\{+\infty\}$, $l_j\le u_j$.

\subsection{The dual simplex method}
\label{sec:dual-simplex}

We fix the notation used below for the dual simplex method with bound constraints \cite{lemke1954,maros2003,koberstein2005}. The index set is partitioned into a basis $B$ ($|B|=m$, $A_B$ nonsingular) and a nonbasis $N$; each nonbasic variable is assigned a status $\eta(j)\in\{L,U,Z\}$, and $x_j=l_j,u_j,0$ according as $\eta(j)=L,U,Z$. The statuses $L,U$ denote placement at a finite lower or upper bound, and status $Z$ keeps a free variable nonbasic at value $0$ \cite{maros2003}. The basic variables are determined by $x_B = A_B^{-1}(b - A_N x_N)$, and $(B,\eta)$ is said to be \emph{primal feasible} if $l_B\le x_B\le u_B$. With the simplex multipliers $y^\top = c_B^\top A_B^{-1}$ and the reduced costs $\bar c_j = c_j - y^\top A_j$, $(B,\eta)$ is said to be \emph{dual feasible} if for all $j\in N$
\[
  \eta(j)=L \ \Longrightarrow\ \bar c_j\ge 0,\qquad
  \eta(j)=U \ \Longrightarrow\ \bar c_j\le 0,\qquad
  \eta(j)=Z \ \Longrightarrow\ \bar c_j=0
\]
holds. A basis that is both primal and dual feasible is an optimal basis.

One iteration starting from a dual feasible basis is as follows.
\begin{enumerate}[label=(\roman*)]
\item \textbf{Leaving row}: choose a row $r$ that violates its bounds, and let $\beta(r)$ be the index of its basic variable. Set $\rho=+1$ if $x_{B,r}<l_{\beta(r)}$ and $\rho=-1$ if $x_{B,r}>u_{\beta(r)}$. Since $x_{B,r}$ is finite, the bound of $\beta(r)$ on the violated side is necessarily finite.
\item \textbf{Pivot row}: let $\alpha^\top := e_r^\top A_B^{-1}A_N$, define the sign $\sigma_j$ to be $+1$ if $\eta(j)=L$, $-1$ if $\eta(j)=U$, and the sign opposite to $\rho\alpha_j$ if $\eta(j)=Z$, and set $\hat\alpha_j:=\sigma_j\alpha_j$, $\hat c_j:=\sigma_j\bar c_j\ (\ge0)$.
\item \textbf{Ratio test}: let $\mathrm{Eligible} := \{\, j\in N : l_j<u_j,\ \rho\,\hat\alpha_j < 0 \,\}$ (for the reason for excluding fixed variables, see Remark~\ref{rem:fixed}). If $\mathrm{Eligible}=\emptyset$, the problem is infeasible and we stop. Otherwise, take $q \in \operatorname*{arg\,min}_{j\in \mathrm{Eligible}} \hat c_j/|\hat\alpha_j|$ as the entering variable. Candidates with status $Z$ have $\hat c_j=0$ and are therefore given top priority with ratio $0$.
\item \textbf{Pivot}: fix $\beta(r)$ at the bound on its violated side, making it nonbasic, and make $q$ basic.
\end{enumerate}
When all rows satisfy their bounds, an optimal basis has been obtained, and under an anti-cycling rule the iterations terminate after finitely many steps \cite{maros2003,koberstein2005}.

\begin{remark}
\label{rem:Z-invariant}
The reduced cost of a nonbasic variable $k$ with status $Z$ remains $0$ throughout the iterations. This is because, in the update $\bar c_k \gets \bar c_k-\theta\,\alpha_k$ ($\theta$ being the dual step length), if $\alpha_k\neq0$ then $k$ can enter with ratio $0$, so that $\theta=0$.
\end{remark}

\section{Construction of an Initial Basis}
\label{sec:difficulty}

A primal feasible initial basis is not trivial in general, and the two-phase method, the big-$M$ method \cite{dantzig1963} and crash procedures \cite{maros2003} are used. On the other hand, for the all-slack basis consisting of slack variables with cost $0$, we have $y=0$ and hence $\bar c_j=c_j$, so dual feasibility is achieved merely by setting $\eta(j)=L$ if $c_j\ge0$ and $\eta(j)=U$ if $c_j<0$. This placement, however, requires the bound at which each variable is placed to be finite, which does not hold in general except in cases such as bounded-variable problems in which all variables have finite bounds on both sides. Methods that remove this obstacle are known as Phase~1 of the dual simplex method; they include minimization of the sum of dual infeasibilities, the cost shifting method, and the artificial bounding method (the big-$M$ method applied to the dual problem), which replaces infinite bounds with a large constant \cite{wolfe1965,maros2003,koberstein2005,koberstein2007,huang2021}.

The numerical difficulty of the big-$M$ method stems from fixing $M$ as a concrete number. For a fixed basic solution, the objective function of the big-$M$ method is an affine function $z^1M+z^0$ of $M$, so comparison for sufficiently large $M$ reduces to lexicographic comparison of the coefficient pairs $(z^1,z^0)$, and the big-$M$ method is equivalent to lexicographic (preemptive) optimization \cite{sherali1983}. Cococcioni and Fiaschi \cite{cococcioni2021} treat $M$ as a non-Archimedean infinite number and realize this idea as an algorithm for the primal big-$M$ method. The point of this paper is to apply this idea to the big-$M$ method applied to the dual problem, that is, to $M$ as a \emph{bound} on variables.

\section{Bounding by $M$ and a Trivial Dual Feasible Basis}
\label{sec:method}

\subsection{Preprocessing: shift of variables with a single infinite bound}
\label{sec:shift-preprocess}

All infinite bounds, including those of free variables, are handled directly by the bounding by $M$ in the next subsection, and we do not perform preprocessing that eliminates free variables from the equality constraints. To simplify the subsequent discussion, we perform the following preprocessing beforehand. For each $j=1,\dots,n$, define the shift amount
\[
  t_j \ :=\
  \begin{cases}
    u_j & (l_j=-\infty,\ u_j<+\infty)\\
    l_j & (l_j>-\infty,\ u_j=+\infty)\\
    0 & (\text{otherwise, i.e., both sides finite or both sides infinite (free variable)})
  \end{cases}
\]
(to avoid confusion with the slack variables $s\in\R^m$, we denote the shift amounts by $t\in\R^n$), and substitute $x_j = x_j' + t_j$. The constraints then become $Ax' = b-At =: b'$, and the bounds become $l_j':=l_j-t_j$, $u_j':=u_j-t_j$. In particular, for each $j$ that had exactly one infinite bound, the bound on the finite side becomes exactly $0$ (for example, if $l_j=-\infty,u_j<\infty$ then $l_j'=-\infty,u_j'=0$; if $l_j>-\infty,u_j=+\infty$ then $l_j'=0,u_j'=+\infty$). For free variables $t_j=0$, and $l_j'=-\infty,u_j'=+\infty$ remain unchanged by the shift. Since the objective function becomes $c^\top x = c^\top x' + c^\top t$, the optimal value of the original problem is the optimal value of the shifted problem plus the constant $c^\top t$.

In what follows we assume that this preprocessing has been done, and we use the symbols $b,l,u$ to mean $b',l',u'$. That is, we may assume that, for variables with a single infinite bound, the bound on the finite side is $0$.

\subsection{Bounding the unbounded variables by $M$}
\label{sec:truncation}

Fix an arbitrary $M>0$. For each $j=1,\dots,n$, define the bounded bounds
\begin{equation}
  \hat l_j :=
  \begin{cases}
    l_j & (l_j > -\infty)\\
    -M & (l_j = -\infty)
  \end{cases}
  ,\qquad
  \hat u_j :=
  \begin{cases}
    u_j & (u_j < +\infty)\\
    M & (u_j = +\infty)
  \end{cases}
  \label{eq:trunc}
\end{equation}
Bounds that were originally finite are used as they are, and only the infinite sides are replaced by $\mp M$. Since all $\hat l_j,\hat u_j$ are finite, the problem
\begin{equation}
  \min\ c^\top x \quad \text{s.t.}\quad Ax=b,\ \hat l\le x\le \hat u
  \label{eq:lpM}
\end{equation}
is a bounded problem to which the simplex method with bound constraints applies directly. Clearly $[\hat l,\hat u]\subseteq[l,u]$, so the feasible region of \eqref{eq:lpM} is a subset of the feasible region of the original problem \eqref{eq:lp}. Moreover, if $M_1\le M_2$ then $[\hat l(M_1),\hat u(M_1)]\subseteq[\hat l(M_2),\hat u(M_2)]$, and the feasible region of \eqref{eq:lpM} is monotonically nondecreasing in $M$.

\subsection{A trivial initial basis via zero-fixed slack variables}
\label{sec:init-basis}

We add to the equality constraints $Ax=b$ slack variables $s\in\R^m$, one for each row, and rewrite them as
\[
  Ax + s = b,\qquad l_s = u_s = 0 \ \ (\text{i.e., } s\equiv 0 \text{ is fixed}).
\]
The extended variable vector is $(x,s)\in\R^{n+m}$, and the extended constraint matrix is $\hat A = [\,A \mid I_m\,]$. The cost of $s$ is $0$.

If we now adopt $B_0:=\{n+1,\dots,n+m\}$ (the index set of the slack variables) as the initial basis, then $A_{B_0}=I_m$, so
\[
  A_{B_0}^{-1} = I_m
\]
and not even the computation of an inverse is required. The statuses of the nonbasic (original) variables are determined solely from the signs of the cost coefficients as
\begin{equation}
  \eta(j) =
  \begin{cases}
    L, \ x_j=\hat l_j & c_j>0,\ \text{or}\ (c_j=0\ \text{and}\ l_j>-\infty)\\
    U, \ x_j=\hat u_j & c_j<0,\ \text{or}\ (c_j=0\ \text{and}\ l_j=-\infty,\ u_j<+\infty)\\
    Z, \ x_j=0 & c_j=0\ \text{and}\ l_j=-\infty,\ u_j=+\infty
  \end{cases}
  \qquad (j=1,\dots,n).
  \label{eq:init-status}
\end{equation}
A variable with cost $0$ satisfies the dual feasibility condition whichever of $L,U,Z$ it is placed at, so we choose a placement whose value does not depend on $M$. That is, if there is a finite side, the variable is placed there, and a free variable is placed at value $0$ with status $Z$. In this way, the only variables whose initial values depend on $M$ are those belonging to the set of variables for which the sign of the cost strictly requires placement on an infinite side,
\[
  S \ :=\ \{\,j : l_j=-\infty,\ c_j>0\,\} \ \cup\ \{\,j : u_j=+\infty,\ c_j<0\,\}.
\]
The values of the basic (slack) variables are then computed directly by
\begin{equation}
  x_{B_0} = s = b - A x_N.
  \label{eq:init-basic}
\end{equation}
In general $s$ does not coincide with $0$, and the value in \eqref{eq:init-basic} may violate the bounds $[0,0]$ of $s$ itself; as we shall see below, however, this poses no problem from the standpoint of dual feasibility.

\begin{proposition}[Dual feasibility of the initial basis]
\label{prop:dual-feasible-init}
The pair $(B_0,\eta)$ defined by \eqref{eq:init-status} is dual feasible for problem~\eqref{eq:lpM}.
\end{proposition}

\begin{proof}
The basic costs are $c_{B_0}=0$ (the cost of $s$ is $0$) and $A_{B_0}=I$, so the simplex multipliers are $y^\top = c_{B_0}^\top A_{B_0}^{-1} = 0$. Hence the reduced cost of a nonbasic original variable $j$ is
\[
  \bar c_j = c_j - y^\top A_j = c_j,
\]
which coincides with the cost itself. By \eqref{eq:init-status}, if $\eta(j)=L$ then $c_j>0$ or ($c_j=0$ and $l_j>-\infty$); in either case $c_j\ge0$, hence $\bar c_j=c_j\ge0$. Similarly, if $\eta(j)=U$ then $c_j<0$ or ($c_j=0$ and $l_j=-\infty$); in either case $c_j\le0$, hence $\bar c_j=c_j\le0$. If $\eta(j)=Z$ then $c_j=0$, so $\bar c_j=0$. Therefore the dual feasibility conditions are satisfied. The reduced costs of the basic variables ($s$) are identically $0$ by definition, and no condition needs to be imposed on them.
\end{proof}

What is essentially used here are two points: (a) since $c_{B_0}=0$ and $A_{B_0}=I$, the simplex multipliers are trivially $y=0$; and (b) the bounds $\hat l_j,\hat u_j$ at which variables are placed are all finite. Point (b) is guaranteed by the bounding by $M$.

\subsection{The algorithm with $M$ fixed}
\label{sec:M-algorithm}

From the above, the algorithm for problem~\eqref{eq:lpM}, bounded with $M$ fixed, is as follows.

\textbf{Step~1 (Construction of the initial basis).} By \eqref{eq:init-status} and \eqref{eq:init-basic}, set the basis $B\gets B_0=\{n+1,\dots,n+m\}$ (the indices of the slack variables), the statuses of the nonbasic original variables to $\eta$, the nonbasic values to $x_N$, and the basic values to $x_B=s=b-Ax_N$. By Proposition~\ref{prop:dual-feasible-init}, this $(B,\eta)$ is dual feasible without any computation.

\textbf{Step~2 (Iteration).} Repeat the following (a)--(d) until termination is determined in (a).

\begin{enumerate}[label=(\alph*)]
\item \textbf{Test of primal feasibility (termination condition).} The values of the basic variables are $x_B=A_B^{-1}(b-A_Nx_N)$, and the nonbasic values are $x_j=\hat l_j,\hat u_j,0$ according as $\eta(j)=L,U,Z$ (in the first pass these coincide with \eqref{eq:init-basic}, and thereafter they may be updated successively by the pivots in (d)). In particular, if originally $l_j=-\infty$ and $\eta(j)=L$ then $x_j=-M$, and if originally $u_j=+\infty$ and $\eta(j)=U$ then $x_j=M$. Let $\beta(i)$ be the index of the basic variable corresponding to row $i$, and define the set of violating rows
\[
  V \ :=\ \bigl\{\, i\in\{1,\dots,m\} \ :\ x_{B,i} < \hat l_{\beta(i)} \ \text{or}\ x_{B,i} > \hat u_{\beta(i)} \,\bigr\}.
\]
If $V=\emptyset$, then $(B,\eta,x)$ is primal feasible for problem~\eqref{eq:lpM}; since dual feasibility holds in Step~1 and is preserved in (b)--(d) (Section~\ref{sec:dual-simplex}), it is an optimal solution of problem~\eqref{eq:lpM}. The algorithm terminates. If $V\neq\emptyset$, proceed to (b).

\item \textbf{Selection of the leaving variable.} First, for each $i\in V$, define the violation
\[
  \Delta_i \ :=\ \max\bigl(\hat l_{\beta(i)}-x_{B,i},\ x_{B,i}-\hat u_{\beta(i)}\bigr) \ (>0).
\]
The leaving row $r\in V$ is chosen by the largest-violation rule (Dantzig's rule) $r\in\operatorname*{arg\,max}_{i\in V}\Delta_i$, or by the steepest-edge rule \cite{forrestgoldfarb1992}, which maximizes $\Delta_i^2/\gamma_i$ normalized by the row weights $\gamma_i:=\sum_{j\in N}\alpha_{ij}^2$ ($\alpha_{ij}:=(A_B^{-1}A_N)_{ij}$), or by the Devex rule \cite{harris1973}, which replaces $\gamma_i$ with approximate weights $\tilde\gamma_i$ that are cheap to update (the $M\to\infty$ version of Dantzig's rule is given in Section~\ref{sec:M-infty}, and the steepest-edge and Devex rules are used as they are for the problem of each phase in the final algorithm of Section~\ref{sec:summary}). The direction of violation $\rho$ is determined with respect to $\hat l,\hat u$ as in (i) of Section~\ref{sec:dual-simplex}.

\item \textbf{Computation of the pivot row and ratio test.} As in (ii) of Section~\ref{sec:dual-simplex}, define $\alpha^\top := e_r^\top A_B^{-1}A_N$, $\sigma_j$, $\hat\alpha_j:=\sigma_j\alpha_j$, $\hat c_j:=\sigma_j\bar c_j\,(\ge0)$, and compute the set of nonbasic variables eligible to enter,
\[
  \mathrm{Eligible} \ :=\ \{\, j\in N : \hat l_j<\hat u_j,\ \rho\,\hat\alpha_j < 0 \,\}
\]
(fixed variables are excluded; see Remark~\ref{rem:fixed} below). If $\mathrm{Eligible}=\emptyset$, problem~\eqref{eq:lpM} is infeasible under the current $M$, and the algorithm terminates. Otherwise, choose the entering variable $q$ by
\[
  q \ \in\ \operatorname*{arg\,min}_{j\in \mathrm{Eligible}} \frac{\hat c_j}{|\hat\alpha_j|}
\]
(candidates with status $Z$ are given top priority with ratio $0$).

\item \textbf{Pivot.} Make $\beta(r)$ nonbasic, fixing it at $\hat l_{\beta(r)}$ if $\rho=+1$ and at $\hat u_{\beta(r)}$ if $\rho=-1$. Make $q$ basic, update the basis $B$, the inverse $A_B^{-1}$, the basic solution $x_B$, the nonbasic values $x_N$ and the nonbasic statuses $\eta$, and return to (a).
\end{enumerate}

\begin{remark}[Exclusion of fixed variables]
\label{rem:fixed}
A nonbasic variable with $\hat l_j=\hat u_j$ (a fixed variable) satisfies dual feasibility without changing its value, merely by switching its status between $L$ and $U$ according to the sign of its reduced cost; hence it may be excluded from the ratio test \cite{koberstein2005}. In particular, since the slack variables $s$ have $l_s=u_s=0$, a slack variable that has once left the basis never returns to it.
\end{remark}

Running this algorithm with a numerical $M$ fixed gives rise to two problems. First, the nonbasic values $\mp M$ enter the computation of $x_B$, the test of violations and the ratio test, causing numerical instability of the same kind as in the classical big-$M$ method. Second, and more fundamentally, unless $M$ is guaranteed to be large enough, the verdict itself may be wrong. Even if the original problem is unbounded, if $M$ is too small then \eqref{eq:lpM} becomes a bounded problem in which the unbounded directions are cut off, and a finite optimal solution is returned. In the next section we avoid these problems by not fixing the value of $M$ and analyzing the behavior as $M\to+\infty$ symbolically.

\section{Symbolic analysis as $M\to+\infty$}
\label{sec:symbolic}

\subsection{Affinity of the basic solution and lexicographic comparison}
\label{sec:M-infty}

As noted in Section~\ref{sec:difficulty}, in the big-$M$ method a comparison for sufficiently large $M$ reduces to a lexicographic comparison of coefficient pairs. Since $M$ in this paper plays the role of a penalty coefficient in the dual problem (Section~\ref{sec:dual-view}), the limit $M\to+\infty$ can likewise be handled as a lexicographic comparison of coefficients.

Steps~2(c)(d) are determined solely by the current basis $B$ and the status $\eta$ of the nonbasic variables, and do not depend on the numerical value of $M$. $M$ enters only through $\hat l,\hat u$, through $x_N,x_B$ determined by them, and through the tests and selections in Steps~2(a)(b). In what follows, we fix the combinatorial state $(B,\eta)$ and carry out the analysis using the fact that these quantities are affine functions of $M$.

\begin{lemma}[Affinity of $x_B$ in $M$]
\label{lem:affine}
Fix the combinatorial state $(B,\eta)$. Each component of the nonbasic values $x_N(M)$ is
\[
  x_j(M) =
  \begin{cases}
    l_j & (\eta(j)=L,\ l_j>-\infty)\\
    -M & (\eta(j)=L,\ l_j=-\infty)\\
    u_j & (\eta(j)=U,\ u_j<+\infty)\\
    M & (\eta(j)=U,\ u_j=+\infty)\\
    0 & (\eta(j)=Z)
  \end{cases}
\]
so we can write $x_N(M) = x_N^0 + Mx_N^1$. Here $x_N^0\in\R^{|N|}$ and $x_N^1\in\{-1,0,1\}^{|N|}$ are constant vectors that depend only on $(B,\eta)$ and not on $M$, with $x_j^1=-1\,(\eta(j)=L,l_j=-\infty)$, $x_j^1=+1\,(\eta(j)=U,u_j=+\infty)$, and $x_j^1=0$ otherwise (including status $Z$). Hence
\[
  x_B(M) \ =\ A_B^{-1}\bigl(b-A_Nx_N(M)\bigr) \ =\ \underbrace{A_B^{-1}(b-A_Nx_N^0)}_{=:x_B^0}\ +\ M\underbrace{\bigl(-A_B^{-1}A_Nx_N^1\bigr)}_{=:x_B^1}
\]
is an affine function of $M$, and the coefficients $x_B^0,x_B^1\in\R^m$ depend only on $(B,\eta)$. That is, setting $x^0:=(x_N^0,x_B^0)$ and $x^1:=(x_N^1,x_B^1)$, we have $x(M)=x^0+Mx^1$.
\end{lemma}

Throughout what follows, for an affine function $f(M)$ of $M$ we denote its \emph{slope} (the coefficient of $M$) by a superscript $1$ and its \emph{intercept} (the constant term) by a superscript $0$, and write $f(M)=f^1M+f^0$. The basic solution $x(M)=x^0+Mx^1$ follows this notation.

Similarly, by \eqref{eq:trunc}, $\hat l_{\beta(i)}(M)$ and $\hat u_{\beta(i)}(M)$ are the affine functions ($l_{\beta(i)}$ or $-M$) and ($u_{\beta(i)}$ or $M$), respectively, so by Lemma~\ref{lem:affine} the violations in Step~2(a),
\[
  v_i^-(M) := \hat l_{\beta(i)}(M) - x_{B,i}(M), \qquad
  v_i^+(M) := x_{B,i}(M) - \hat u_{\beta(i)}(M),
\]
are also affine functions of $M$. Following the notation above, we write them as $v_i^-(M)=v_i^{-,1}M+v_i^{-,0}$ and $v_i^+(M)=v_i^{+,1}M+v_i^{+,0}$ (where $v_i^{\pm,1},v_i^{\pm,0}\in\R$ are constants depending only on $(B,\eta)$).

\begin{definition}[Lexicographic order on affine functions]
For two affine functions $f(M)=f^1M+f^0$ and $h(M)=h^1M+h^0$, define
\[
  f \ \succ\ h \quad :\Longleftrightarrow\quad f^1>h^1,\ \ \text{or}\ \ (f^1=h^1\ \text{and}\ f^0>h^0)
\]
(that is, we compare the coefficient vectors $(f^1,f^0),(h^1,h^0)$ in lexicographic order).
\end{definition}

If $f\succ h$, then there exists $M^\dagger(f,h)>0$ such that $f(M)>h(M)$ for all $M\ge M^\dagger(f,h)$ (if $f^1\neq h^1$ one may take $M^\dagger=|f^0-h^0|/|f^1-h^1|$, and if $f^1=h^1,\,f^0>h^0$ the inequality holds for every $M>0$). The case $h\succ f$ is analogous, and when $(f^1,f^0)=(h^1,h^0)$ we have $f\equiv h$. This $\succ$ gives a total order on any finite set of affine functions, and if $M^\dagger$ denotes the maximum of $M^\dagger(f,h)$ over all pairs among them, then as long as $M\ge M^\dagger$, the outcomes of all numerical comparisons coincide with the comparisons by $\succ$.

Using this, Steps~2(a)(b) can be rewritten as follows without fixing a specific value of $M$.

\begin{enumerate}
\item[(a$'$)] \textbf{Primal feasibility test (termination condition, $M\to\infty$ version).} Let $V_\infty := \{\, i\in\{1,\dots,m\} : v_i^- \succ 0 \ \text{or}\ v_i^+ \succ 0 \,\}$ (where $0$ denotes the affine function that is identically $0$, i.e., the coefficient pair $(0,0)$). If $V_\infty=\emptyset$, primal feasibility holds for all $M\ge M^\dagger$, and the algorithm terminates.
\item[(b$'$)] \textbf{Selection of the leaving variable (Dantzig rule, $M\to\infty$ version).} For each $i\in V_\infty$, let $w_i$ be the larger of $v_i^-,v_i^+$ with respect to $\succ$. Choose the leaving row $r$ by
\[
  r \ \in\ \operatorname*{arg\,max}_{i\in V_\infty}{}^{\succ}\, w_i
\]
(as a maximal element with respect to the lexicographic order $\succ$). Set $\rho=+1$ if $w_r=v_r^-$ and $\rho=-1$ if $w_r=v_r^+$ (ties, when $\arg\max^\succ$ contains several rows, are broken deterministically by a rule in the spirit of Bland's rule, such as giving priority to the smaller index).
\end{enumerate}

\begin{proposition}[Agreement with numerical computation for sufficiently large $M$]
\label{prop:consistency}
For a fixed combinatorial state $(B,\eta)$, there exists $M^\dagger>0$ (depending only on that state) such that the tests and selections by (a$'$)(b$'$) coincide with the tests and selections by the numerical computation in (a)(b) for all $M\ge M^\dagger$.
\end{proposition}

\begin{proof}
The membership test for each element of $V_\infty$ and the choice of $r$ by $\arg\max^\succ$ are determined solely by comparisons with respect to $\succ$ among the finitely many affine functions $v_i^-$, $v_i^+$ and $0$. As noted above, beyond some common $M^\dagger$ these comparisons coincide with the numerical comparisons $v_i^\pm(M)>0$ and $\max_i(\cdot)$.
\end{proof}

Since Steps~2(c)(d) do not depend on $M$, (a$'$)(b$'$) uniquely determine the next state $(B',\eta')$ (without fixing $M$ as a number), and repeating this yields a sequence of combinatorial state transitions $(B^{(0)},\eta^{(0)}) \to (B^{(1)},\eta^{(1)}) \to \cdots$. Under an anti-cycling rule (Section~\ref{sec:anticycling}) this sequence terminates after finitely many steps, and if we redefine $M^\dagger$ as the maximum of the $M^\dagger$ over all iterations, then for every $M\ge M^\dagger$ the sequence of iterates obtained by numerically computing Step~2 coincides exactly with this combinatorial transition sequence. In other words, the behavior as $M\to+\infty$ can be computed solely by lexicographic comparisons of coefficient pairs. The final state $(B^\ast,\eta^\ast)$ of this transition sequence is itself an optimal basis of \eqref{eq:lpM} for sufficiently large $M$. Hereafter, we call this iteration in which $M$ is treated symbolically (Steps~2(a$'$)(b$'$)(c)(d)) the \emph{symbolic artificial bounding method}.

\subsection{Classification into the three cases}
\label{sec:trichotomy}
Suppose the state transition sequence has reached the final state $(B^\ast,\eta^\ast)$ (a state in which $V_\infty=\emptyset$ in (a$'$)). Using $x(M)=x^0+Mx^1$ from Lemma~\ref{lem:affine}, the optimal value of problem~\eqref{eq:lpM} is again expressed as an affine function of $M$:
\[
  z(M) \ =\ c_B^\top x_B(M) + c_N^\top x_N(M)
     \ =\ \underbrace{c^\top x^0}_{=:z^0} \ +\ M\ \underbrace{c^\top x^1}_{=:z^1} .
\]
Here $z^0,z^1\in\R$ depend only on $(B^\ast,\eta^\ast)$, and for $M\ge M^\dagger$ this is exactly the optimal value of problem~\eqref{eq:lpM}. Since the feasible region of \eqref{eq:lpM} is monotonically nondecreasing in $M$ (Section~\ref{sec:truncation}), $z(M)$ is monotonically nonincreasing, and hence $z^1\le0$ always holds.

\begin{lemma}[Union of feasible regions]
\label{lem:union}
Let $F:=\{x:Ax=b,\ l\le x\le u\}$ (the feasible region of the original problem~\eqref{eq:lp}). Then
\[
  F \ =\ \bigcup_{M>0} F_M, \qquad F_M := \{x:Ax=b,\ \hat l(M)\le x\le \hat u(M)\}\ (\text{i.e., the feasible region of \eqref{eq:lpM}})
\]
holds.
\end{lemma}

\begin{proof}
($\supseteq$): For any $M>0$ we have $[\hat l(M),\hat u(M)]\subseteq[l,u]$, so $F_M\subseteq F$.

($\subseteq$): Take any $x\in F$. Since $x$ is a finite real vector (not extended-real-valued), $\|x\|_\infty<\infty$. Taking $M\ge\|x\|_\infty$, for coordinates where $x_j$ originally has finite bounds, $\hat l_j(M)=l_j\le x_j\le u_j=\hat u_j(M)$ holds as is, and for coordinates that were originally unbounded, $|x_j|\le M$ gives $-M\le x_j\le M$. Hence $x\in F_M$.
\end{proof}

From the sign of the slope $z^1$ and the termination outcome, the three cases of the original problem can be distinguished exactly and exhaustively.

\begin{proposition}[Classification into the three cases]
\label{prop:trichotomy}
Suppose the combinatorial state transition sequence reaches the final state $(B^\ast,\eta^\ast)$ beyond some $M^\dagger>0$, where (a$'$)(b$'$)(c)(d) terminates. Then the following hold.
\begin{enumerate}[label=(\roman*)]
\item The termination outcome is \textsc{Infeasible} $\iff$ $F=\emptyset$.
\item The termination outcome is \textsc{Optimal} and $z^1<0$ $\iff$ $F\neq\emptyset$ and $\inf_{x\in F}c^\top x=-\infty$ (the original problem is feasible and unbounded).
\item The termination outcome is \textsc{Optimal} and $z^1=0$ $\iff$ $F\neq\emptyset$ and $\inf_{x\in F}c^\top x$ is finite, in which case $\inf_{x\in F}c^\top x=z^0$.
\end{enumerate}
\end{proposition}

\begin{proof}
By Proposition~\ref{prop:consistency} and the discussion following it, the termination outcome is the same for all $M\ge M^\dagger$, and either $F_M$ is empty for all $M\ge M^\dagger$, or it is nonempty for all $M\ge M^\dagger$ with minimum value $z(M)=z^0+z^1M$. Moreover, $M_1\le M_2$ implies $F_{M_1}\subseteq F_{M_2}$ (monotonicity).

(i): ($\Leftarrow$) If $F=\emptyset$, then by Lemma~\ref{lem:union}, $F_M\subseteq F=\emptyset$ for every $M$. In particular, \eqref{eq:lpM} is infeasible for $M\ge M^\dagger$, and the termination outcome is \textsc{Infeasible}.
($\Rightarrow$) If the termination outcome is \textsc{Infeasible}, then $F_{M^\dagger}=\emptyset$. By monotonicity, $F_M\subseteq F_{M^\dagger}=\emptyset$ also for $M<M^\dagger$. Hence, by Lemma~\ref{lem:union}, $F=\bigcup_M F_M=\emptyset$.

(ii)(iii): In what follows, by (i) we may assume $F\neq\emptyset$, so the termination outcome is \textsc{Optimal}. By Lemma~\ref{lem:union} and monotonicity, for any $x\in F$, setting $M_x:=\max(M^\dagger,\|x\|_\infty)$, we have $x\in F_{M_x}$ (by the same argument as in the proof of Lemma~\ref{lem:union}), and therefore
\[
  c^\top x \ \ge\ \min_{F_{M_x}} c^\top x \ =\ z(M_x) \ =\ z^0+z^1M_x .
  \tag{$\ast$}
\]

($z^1<0 \Rightarrow$ unbounded): Since $F_M\subseteq F$, $\inf_{x\in F}c^\top x\le z(M)=z^0+z^1M\to-\infty\ (M\to\infty)$.

(Unbounded $\Rightarrow z^1<0$): We prove the contrapositive. Suppose $z^1=0$. If $F\neq\emptyset$ and the problem is unbounded, then for any $N>0$ there exists $x^{(N)}\in F$ with $c^\top x^{(N)}<-N$. Applying $(\ast)$ to $x=x^{(N)}$ gives $c^\top x^{(N)}\ge z^0+z^1M_{x^{(N)}}=z^0$ (independent of the value of $M_{x^{(N)}}$ since $z^1=0$). Hence $z^0\le c^\top x^{(N)}<-N$. Since $N$ is arbitrary, this gives $z^0=-\infty$, contradicting the fact that $z^0$ is defined as a finite real number by Lemma~\ref{lem:affine}. Hence, if the problem is unbounded, then $z^1\neq0$. Combined with $z^1\le0$ shown earlier, $z^1<0$.

From the above, when $F\neq\emptyset$, ``$z^1<0\iff$ unbounded'' holds, and therefore (by the dichotomy given by $z^1\le0$) ``$z^1=0\iff$ a finite optimal value exists'' follows.

Finally, we show that $\inf_{x\in F}c^\top x=z^0$ in the case $z^1=0$. For any $x\in F$, using $z^1=0$ in $(\ast)$ gives $c^\top x\ge z^0$. On the other hand, fixing an arbitrary $M\ge M^\dagger$, since the final state is primal feasible for $M\ge M^\dagger$, the basic solution $(x_B(M),x_N(M))$ at this $M$ belongs to $F_M$, and $F_M\subseteq F$ by the ($\supseteq$) direction of the proof of Lemma~\ref{lem:union}. Moreover, $c_B^\top x_B(M)+c_N^\top x_N(M)=z(M)=z^0+z^1M=z^0$ (the last equality by $z^1=0$). Hence $z^0$ is a lower bound actually attained on $F$, and $z^0=\min_{x\in F}c^\top x$ follows.
\end{proof}

By Proposition~\ref{prop:trichotomy}, the distinction between an insufficient $M$ and unboundedness, which was problematic in the artificial bounding method with a numerical $M$ (Section~\ref{sec:intro}), can be made exactly, without restarting, from the sign of the slope $z^1$ at termination alone.

In case (iii) of Proposition~\ref{prop:trichotomy}, the optimal \emph{basis} $(B^\ast,\eta^\ast)$ and the optimal value $z^0$ are obtained without fixing $M$, but extracting a concrete optimal solution vector requires the cleanup described below.

\subsection{Cleanup of nonbasic variables remaining at artificial bounds}
\label{sec:cleanup}
Consider the variables $j$ that remain nonbasic in the final state $(B^\ast,\eta^\ast)$ with $x_j^1\neq0$ (i.e., placed on the side that was infinite, so that $x_j(M)=\mp M$). Such a $j$ is either a variable with a single infinite bound (whose bound on the finite side is exactly $0$ by preprocessing) or a free variable.

By the standard simplex argument, rewriting $z(M)=c_B^\top x_B(M)+c_N^\top x_N(M)$ using $x_B(M)=A_B^{-1}(b-A_Nx_N(M))$ gives
\[
  z(M) \ =\ c_B^\top A_B^{-1}b \ +\ \sum_{j\in N} \bar c_j\, x_j(M) .
\]
Extracting only the terms among the nonbasic values $x_j(M)$ that depend on $M$ (those $j$ placed at artificial bounds), the slope of $z(M)=z^0+z^1M$ is expressed as
\[
  z^1 \ =\ \sum_{\substack{j\in N,\ \eta(j)=L\\ l_j=-\infty}} (-\bar c_j) \ +\ \sum_{\substack{j\in N,\ \eta(j)=U\\ u_j=+\infty}} \bar c_j
\]
(this expression equals $z^1=c^\top x^1$ from Lemma~\ref{lem:affine} rewritten in terms of reduced costs). By dual feasibility, each term with $\eta(j)=L$ has $\bar c_j\ge0$ and hence $-\bar c_j\le0$, and each term with $\eta(j)=U$ has $\bar c_j\le0$; thus $z^1$ is a \emph{sum of nonpositive terms}, which also gives an alternative proof of $z^1\le0$ shown earlier. In particular,
\[
  z^1=0 \quad\Longleftrightarrow\quad \bar c_j=0\ \text{for all nonbasic variables $j$ with}\ x_j^1\neq0
\]
holds. This is nothing but the standard optimality condition that the reduced cost of a nonbasic variable placed in an unbounded direction must be $0$.

However, when there exist nonbasic variables with $x_j^1\neq0$, the point $x^0=(x_N^0,x_B^0)$ obtained by formally substituting $M=0$ is in general not a feasible solution of the original problem. Although the nonbasic values $x_N^0$ lie on the true bounds by preprocessing, primal feasibility of the basic values $x_B(M)=x_B^0+Mx_B^1$ is guaranteed only for $M\ge M^\dagger$, and extrapolating down to $M=0$ may violate the bounds. The following lemma resolves this issue by exploiting the property $\bar c_j=0$.

\begin{lemma}[Removal of nonbasic variables at artificial bounds by a primal ratio test]
\label{lem:cleanup}
Suppose $z^1=0$. Assume that the combinatorial state $(B,\eta)$ is dual feasible for the original problem~\eqref{eq:lp}, and that each component of its basic solution $x(M)$ is an affine function of $M$ satisfying the bounds of the original problem $l\le x(M)\le u$ for all sufficiently large $M$ (the final state $(B^\ast,\eta^\ast)$ satisfies this, since $F_M\subseteq F$ and dual feasibility does not depend on $M$). Let $j\in N$ be a nonbasic variable with $x_j^1\neq0$ (so $\bar c_j=0$ by $z^1=0$). By preprocessing, the current value of $j$ is $x_j(M)=x_j^1M$, and if $j$ has a single infinite bound, its bound on the finite side is exactly $0$. If $x_j$ is moved from $x_j^1M$ toward $0$ by a step $t\in[0,M]$ ($x_j\gets x_j^1(M-t)$), the basic variables change as
\[
  x_B \ \gets\ x_B(M) + x_j^1\,t\,\alpha_{\cdot j}, \qquad \alpha_{\cdot j}:=A_B^{-1}A_j .
\]
For each row $i$, define the step at which $x_{B,i}$ reaches a \emph{genuinely finite} bound of $\beta(i)$ as
\[
  t_i :=
  \begin{cases}
    \dfrac{u_{\beta(i)}-x_{B,i}(M)}{x_j^1\alpha_{ij}} & (x_j^1\alpha_{ij}>0,\ u_{\beta(i)}<+\infty)\\[2ex]
    \dfrac{x_{B,i}(M)-l_{\beta(i)}}{-x_j^1\alpha_{ij}} & (x_j^1\alpha_{ij}<0,\ l_{\beta(i)}>-\infty)
  \end{cases}
\]
(the other rows do not block the move; the artificial bounds $\pm M$ do not exist in the original problem and are therefore not taken into account), and let $t^\ast$ be the minimal element of $\{M\}\cup\{t_i\}$ with respect to $\succ$ (all of these are affine functions of $M$; ties are broken in favor of $M$, and among the $t_i$ in favor of the smaller index). Then:
\begin{enumerate}[label=(\Alph*)]
\item If $t^\ast=M$, set $x_j$ to $0$. That is, if $j$ has a single infinite bound, place it on its finite side (status $L$ or $U$); if it is a free variable, place it in status $Z$. The basis is not changed.
\item If $t^\ast=t_r\prec M$, perform a pivot with $j$ as the entering variable and $\beta(r)$ as the leaving variable, and fix $\beta(r)$ at the finite bound it has reached (status $U$ for an upper bound, status $L$ for a lower bound).
\end{enumerate}
Then the following hold.
\begin{enumerate}[label=(\roman*)]
\item The objective value does not change.
\item Dual feasibility for the original problem~\eqref{eq:lp} is preserved.
\item The updated basic solution is also an affine function of $M$ and satisfies $l\le x(M)\le u$ for all sufficiently large $M$.
\item The number of nonbasic variables with $x_j^1\neq0$ decreases by exactly $1$.
\end{enumerate}
\end{lemma}

\begin{proof}
(i): The change in objective value due to the move is $c_j(-x_j^1t)+c_B^\top(x_j^1t\,\alpha_{\cdot j})=-x_j^1t\,(c_j-c_B^\top A_B^{-1}A_j)=-x_j^1t\,\bar c_j=0$.

(ii): In (A), the basis does not change, so the reduced costs are unchanged; only the status of $j$ changes, and $\bar c_j=0$ satisfies the condition for each of the statuses $L,U,Z$. In (B), by the tableau update formula, the reduced cost of every other nonbasic variable $k$ is unchanged, $\bar c_k^{\mathrm{new}}=\bar c_k-(\bar c_j/\alpha_{rj})\alpha_{rk}=\bar c_k$, and the reduced cost of the leaving variable becomes $\bar c_{\beta(r)}^{\mathrm{new}}=-\bar c_j/\alpha_{rj}=0$, so the sign condition is satisfied regardless of which bound $\beta(r)$ is placed at.

(iii): Since $t^\ast$ is an affine function of $M$, each updated component is also an affine function. By assumption, $x(M)$ satisfies the genuinely finite bounds for sufficiently large $M$, so each $t_i\succeq0$, and therefore $0\preceq t^\ast\preceq M$. The nonbasic variable $j$ stays between $x_j^1M$ and $0$, which is within the bounds of $j$. For basic variables $i$ that may block the move, $t^\ast\preceq t_i$ ensures that the slack to the bound remains $\succeq0$; the other basic variables either move toward an infinite side or do not move. The variable $\beta(r)$ leaving in (B) lies exactly on a finite bound. The other nonbasic variables do not change. Here $\succeq0$ coincides with numerical nonnegativity for sufficiently large $M$.

(iv): In (A), the nonbasic value of $j$ becomes $0$, so $x_j^1=0$. In (B), $j$ enters the basis, and $\beta(r)$, which becomes nonbasic in its place, is placed at a finite bound ($x_{\beta(r)}^1=0$). The statuses of the other nonbasic variables do not change.
\end{proof}

We apply Lemma~\ref{lem:cleanup} repeatedly until no nonbasic variables with $x_j^1\neq0$ remain. By (iv), if $K$ denotes the number of such variables in the final state $(B^\ast,\eta^\ast)$, this terminates after exactly $K$ applications. In the resulting basis $(\tilde B,\tilde\eta)$ we have $x_N^1=0$, so the nonbasic values $x_N=x_N^0$ do not depend on $M$, and hence the basic values $x_B=A_{\tilde B}^{-1}(b-A_Nx_N^0)$ are also constants independent of $M$. By (iii), this constant vector $x^{\ast\ast}$ satisfies $l\le x^{\ast\ast}\le u$ for sufficiently large $M$ (and hence regardless of $M$), and is primal feasible for the original problem~\eqref{eq:lp}. Moreover, by (ii) it is dual feasible, so $x^{\ast\ast}$ is an optimal \emph{basic} solution of the original problem, and by (i) $c^\top x^{\ast\ast}=z(M)=z^0$ (consistent with Proposition~\ref{prop:trichotomy}(iii)).

\begin{remark}[Use of status $Z$, and the case of a zero column]
\label{rem:state-F}
Status $Z$ appears only in the initial placement of free variables with cost $0$ (\eqref{eq:init-status}) and when a free variable $j$ is set to value $0$ in (A) of Lemma~\ref{lem:cleanup}. Since variables leaving the basis in the main iterations are placed at $L$ or $U$, no new variable with status $Z$ arises during the main iterations. Also, since $\beta(r)$ leaving in (B) of Lemma~\ref{lem:cleanup} is a variable that has reached a genuinely finite bound, no free variable leaves the basis during cleanup. Furthermore, when $\alpha_{\cdot j}=A_B^{-1}A_j=0$, no $t_i$ is defined, so case (A) applies automatically, and it suffices to move $j$ to its finite side (or to status $Z$) without changing $x_B$.
\end{remark}

\begin{remark}[Comparison with cleanup that pivots on an arbitrary row]
Statements (i)(ii)(iv) of Lemma~\ref{lem:cleanup} also hold for a degenerate pivot that brings $j$ with $x_j^1\neq0$ into the basis on an \emph{arbitrary} row $r$ with $\alpha_{rj}\neq0$ and fixes $\beta(r)$ on its finite side. In this case, however, primal feasibility is lost because the value of $\beta(r)$ is forcibly moved to a bound, and the dual simplex method must be restarted after reaching $x_N^1=0$. The number of iterations of this restart cannot be estimated in advance. Lemma~\ref{lem:cleanup} preserves primal feasibility by choosing the leaving row via a ratio test, thereby making this restart unnecessary. Alternatively, an optimal solution can also be obtained by substituting a concrete numerical value of $M$ (the smallest value that yields primal feasibility) into the basic solution $x(M)$ of the final state, but the resulting point is in general not a basic solution, and if some row has a small slope, $M$ becomes extremely large and causes numerical errors. The procedure of Lemma~\ref{lem:cleanup} treats $M$ symbolically throughout and finally yields a basic solution free of $M$, so this problem does not arise either.
\end{remark}

\begin{remark}[Relation to crossover procedures]
\label{rem:crossover}
\emph{Crossover}~\cite{megiddo1991,bixby1994,andersen1996} is a procedure that, starting from a solution in the relative interior of the optimal face obtained by an interior-point method, constructs an optimal basic solution by moving variables lying strictly between their bounds, one at a time, either to a bound or into the basis, while preserving the objective value; the stage that moves the primal variables is called the primal push, and the stage that moves the dual variables is called the dual push.

The operation of Lemma~\ref{lem:cleanup} has the same structure as the primal push in this crossover procedure. Indeed, the basic solution of the final state $(B^\ast,\eta^\ast)$ is optimal for the original problem, but since the nonbasic variables with $x_j^1\neq0$ do not lie on bounds of the original problem, it is not a basic solution (vertex) of the original problem. Such a nonbasic variable $j$ with $x_j^1\neq0$ is, from the viewpoint of the original problem, a nonbasic variable not placed at any finite bound (a so-called superbasic variable), and by $z^1=0$ its reduced cost is $0$. The primal push moves such variables one at a time toward a bound, and if a basic variable reaches a bound first along the way, performs a pivot exchanging the two. The properties that the objective value and primal and dual feasibility are preserved because the reduced cost is $0$, and that each operation decreases the number of such variables by exactly one (Lemma~\ref{lem:cleanup}(i)--(iv)), are also shared. Megiddo~\cite{megiddo1991} showed by an argument of this kind that an optimal basis can be obtained from a pair of primal and dual optimal solutions in at most $n$ pivots. Lemma~\ref{lem:cleanup} applies this primal push to the symbolic setting in which the steps and the basic solution are affine functions of $M$, performing comparisons in the lexicographic order $\succ$. Note that, under the duality of Section~\ref{sec:dual-view}, this corresponds on the dual side to driving out artificial variables that remained in the basis at value $0$ while preserving optimality, i.e., to the dual push of crossover.
\end{remark}

\subsection{Non-return to infinite bounds and interpretation of the method}
\label{sec:no-return}

Once $x_N^1=0$ (the state in which all nonbasic variables are placed at genuinely finite bounds) is attained in the course of the iterations, infinite bounds never reappear as nonbasic variables in subsequent iterations. The crux of this property is the fact that $x_{B,r}$, being a finite real number, can never ``violate'' an infinite bound.

\begin{proposition}[Non-return of nonbasic variables to infinite bounds]
\label{prop:no-return}
Consider the combinatorial state transition sequence $(B^{(0)},\eta^{(0)})\to(B^{(1)},\eta^{(1)})\to\cdots$ generated by the symbolic artificial bounding method (Steps~2(a$'$)(b$'$)(c)(d)). If for some $k_0\ge0$ all nonbasic variables of $(B^{(k_0)},\eta^{(k_0)})$ have $x_j^1=0$ (i.e., are placed at genuinely finite bounds), then the same holds for all $k\ge k_0$.
\end{proposition}

\begin{proof}
We proceed by induction on $k\ge k_0$. For $k=k_0$ the claim holds by assumption. Assume that for some $k\ge k_0$ all nonbasic variables of $(B^{(k)},\eta^{(k)})$ have $x_j^1=0$, and consider one iteration $(B^{(k)},\eta^{(k)})\to(B^{(k+1)},\eta^{(k+1)})$.

First, since $x_N^1=0$ in Lemma~\ref{lem:affine}, we have $x_N(M)\equiv x_N^0$, and therefore $x_B^1=-A_B^{-1}A_Nx_N^1=0$, i.e., $x_B(M)\equiv x_B^0$ is a constant independent of $M$.

For each row $i$, if $\beta(i)$ has an infinite side (e.g., $l_{\beta(i)}=-\infty$), the corresponding violation $v_i^-(M)=\hat l_{\beta(i)}(M)-x_{B,i}(M)=(-M)-x_{B,i}^0$ has slope $-1$ and diverges to $-\infty$ as $M\to\infty$, so $v_i^-\succ0$ never holds for $M\ge M^\dagger$, whatever $M^\dagger$ is. On the other hand, the violation corresponding to a genuinely finite side of $\beta(i)$, together with the fact that $x_{B,i}(M)=x_{B,i}^0$ is constant, is likewise a constant independent of $M$ (the case $u_{\beta(i)}=+\infty$ is symmetric). Therefore, the violation $w_r(M)=w_r^1M+w_r^0$ of the leaving row $r$ is a constant (an ordinary real number) with slope $w_r^1=0$.

(i) \textbf{The leaving variable is fixed at a genuinely finite bound}: Since $w_r(M)$ is constant, $r$ is necessarily selected because of a violation on a genuinely finite side of $\beta(r)$ (the violation on an infinite side diverges with slope $-1$ as seen above and cannot be $\succ0$). By the pivot in Step~2(d), $\beta(r)$ is fixed at that genuinely finite bound and becomes nonbasic.

(ii) \textbf{Conclusion}: The entering variable $q$, which becomes basic in the pivot, is removed from the nonbasic set; by (i), $\beta(r)$, which becomes nonbasic in the pivot, is fixed at a genuinely finite bound; and the other nonbasic variables do not change status in this iteration. Therefore all nonbasic variables of $(B^{(k+1)},\eta^{(k+1)})$ also have $x_j^1=0$.
\end{proof}

\begin{remark}[$x_N^1=0$ as an absorbing boundary and implementation consequences]
\label{rem:absorbing}
From the proof of Proposition~\ref{prop:no-return}, once $x_N^1=0$ is reached, all violations become constants independent of $M$, and the tests in Step~2 (primal feasibility test, selection of the leaving row, ratio test) coincide exactly with those of the dual simplex method of Section~\ref{sec:dual-simplex} with respect to the true bounds $l,u$. That is, reaching $x_N^1=0$ functions as an absorbing boundary that makes the subsequent behavior coincide with the classical dual simplex method.
\end{remark}

\begin{remark}[On iterations after reaching $z^1=0$]
\label{rem:z1zero}
The slope $z^1=\sum_{j\in N}x_j^1\bar c_j$ of $z(M)$ in the current state has only nonpositive terms by dual feasibility, and $-z^1$ corresponds to the dual infeasibility with respect to the original problem~\eqref{eq:lp} (the sum of the absolute values of the reduced costs of the nonbasic variables placed on infinite sides). In the initial state, $z^1=-\sum_{j\in S}|c_j|$. Since each iteration of the dual simplex method does not decrease the objective value for all sufficiently large $M$, $(z^1,z^0)$ is lexicographically nondecreasing; in particular, $z^1$ is monotonically nondecreasing, and once it reaches $0$ it remains $0$ thereafter. In this sense, the symbolic artificial bounding method roughly corresponds to dual Phase~1 while $z^1<0$ and to Phase~2 after reaching $z^1=0$, and if $S=\emptyset$ it is already in the latter at the initial state. However, the precise division into phases is given not by the sign of $z^1$ but by Phase A and Phase B of Section~\ref{sec:two-phase} (Remark~\ref{rem:not-delta}), and the sequence of iterates coincides exactly with the classical dual simplex method only after $x_N^1=0$ is reached (Remark~\ref{rem:absorbing}).
\end{remark}

\begin{remark}[Interpretation as a guaranteed crash procedure]
Since the dual feasibility condition is determined solely by the signs of the reduced costs and does not depend on $M$, the iterations of Step~2 are dual feasible for the original problem~\eqref{eq:lp} as well, consistently from the start. What Step~2 actually resolves is merely the single point that some nonbasic variables are placed at the artificial values $\pm M$ ($x_N^1\neq0$). From this viewpoint, Step~2 (including the cleanup of Lemma~\ref{lem:cleanup}) can be regarded as a \emph{crash procedure} for obtaining a starting basis that is dual feasible for the original problem and in which all nonbasic variables are placed at finite bounds. Whereas the classical crash procedures of Section~\ref{sec:difficulty} are heuristic and carry no feasibility guarantee, this procedure, under an anti-cycling rule (Section~\ref{sec:anticycling}), \emph{always} reaches $x_N^1=0$ after finitely many iterations (or detects unboundedness or infeasibility), and by Proposition~\ref{prop:no-return} never moves back once it has been reached.
\end{remark}

\begin{remark}[Comparison with the cost shifting method]
A technique widely used at the start of the dual simplex method is the \emph{cost shifting method}~\cite{maros2003}. It is a two-phase procedure that temporarily replaces the costs $c_j$ of nonbasic variables violating the dual feasibility conditions by artificial values so as to make the basis \emph{apparently} dual feasible, runs the dual simplex method under this artificial objective function, and then restores dual feasibility by additional pivots while returning the costs to their true values.

From the dual side, both the cost shifting method and our bounding by $M$ are operations that temporarily relax the violated dual constraints $\bar c_j\le0$ or $\bar c_j\ge0$. The difference is that the cost shifting method relaxes them \emph{without penalty} by moving the right-hand side $c_j$ of the dual constraint itself and later restores it, whereas our bounding relaxes them \emph{with a penalty of coefficient $M$}, keeping the violation as an artificial variable (see Section~\ref{sec:dual-view} below). From the primal side, our approach does not artificialize the objective function and uses the true $c^\top x$ from beginning to end. Because the penalty incorporates the elimination of dual infeasibility into the objective function, advantages arise in the following two respects. First, whereas in the cost shifting method the number of additional pivots in the stage that undoes the shift cannot in general be characterized in advance, our cleanup is proved to finish in at most $K$ pivots without changing the objective value and while preserving primal and dual feasibility (Lemma~\ref{lem:cleanup}). Second, whereas in the cost shifting method the variables to be shifted and the shift amounts are chosen heuristically, in our approach the set $S$ of variables requiring placement on an infinite side is uniquely determined by the data, and moreover an exact criterion for classifying the original problem into the three cases (Proposition~\ref{prop:trichotomy}) is obtained within the same framework.
\end{remark}

\subsection{Anti-cycling: extending Bland's rule to the lexicographic order}
\label{sec:anticycling}

So far, in Steps~2(a$'$)(b$'$) as well as in the ratio test, we have only stated that ties among candidates are broken deterministically by a rule such as ``give priority to the smaller index,'' with the caveat that this alone cannot in general exclude cycling (repeatedly visiting the same basis). The standard ways to prevent cycling in the classical dual simplex method with bound constraints are to use Bland's rule (the smallest-index rule)~\cite{bland1977} itself as the primary selection rule, or to use the \emph{lexicographic rule}~\cite{dantzig1955}, which is compatible with any primary selection rule such as the Dantzig rule, the steepest-edge rule, or the Devex rule. The lexicographic rule in the dual simplex method resolves ties in the ratio test (Step~2(c)) by lexicographically comparing the reduced costs obtained when the costs $c$ are perturbed by infinitesimals $\epsilon=(\epsilon_1,\epsilon_2,\dots)$, $\epsilon_1\gg\epsilon_2\gg\cdots\gg0$. Since the cause of cycling is dual degeneracy (ties with ratio $0$ in the ratio test), perturbation is needed only on the ratio-test side, and ties in the selection of the leaving row (Step~2(b)) may be broken by any deterministic rule.

In our setting, this lexicographic rule does not interfere with the lexicographic order $\succ$ with respect to $M$. The ratios $\hat c_j/|\hat\alpha_j|$ compared in the ratio test do not depend on $M$, and neither do the perturbed reduced costs, so ties in the ratio test can be resolved solely by lexicographic comparison with respect to $\epsilon$. On the other hand, the tests in which $M$ appears (comparisons of violations in (a$'$)(b$'$)) are carried out solely with $\succ$ with respect to $M$, and ties in the leaving row may be broken by a rule such as the smallest index. That is, there is no need to prepare a lexicographic order concatenating $M$ and $\epsilon$.

\begin{proposition}[Anti-cycling]
Suppose ties in the ratio test are resolved by the lexicographic rule based on the cost perturbation above, and all other ties are resolved by an arbitrary deterministic rule. Then the combinatorial iterations of the symbolic artificial bounding method never visit the same state $(B,\eta)$ twice, and terminate after finitely many iterations.
\end{proposition}

\begin{proof}
We first handle variables with status $Z$. The costs of variables with status $Z$ are not perturbed. When a variable with status $Z$ is eligible to enter, it becomes basic with dual step length $\theta=0$ (Remark~\ref{rem:Z-invariant}), so the reduced costs are unchanged even after perturbation, and the reduced costs of variables with status $Z$ remain exactly $0$ in the perturbed problem as well. Each such pivot decreases the number of nonbasic variables with status $Z$ by one, and no new variable with status $Z$ arises during the main iterations (Remark~\ref{rem:state-F}). Hence pivots of this kind occur at most as many times as the initial number of variables with status $Z$. In what follows, we consider the other pivots.

In the perturbed problem, the reduced costs of all nonbasic variables except fixed variables and variables with status $Z$ are lexicographically nonzero with respect to $\epsilon$, and it is a standard fact of the classical lexicographic dual simplex method that this property is preserved throughout the iterations. Therefore the dual step length $\theta_q$ in each iteration is lexicographically positive with respect to $\epsilon$. On the other hand, the violation of the leaving row satisfies $w_r\succ0$, and the change in the objective value is $\theta_q\,w_r(M)$, so for sufficiently large $M$ and sufficiently small $\epsilon$ the objective value strictly increases. Since there are finitely many states $(B,\eta)$ and the objective value is determined by the state (and since the states are finite in number, the threshold on $M$ can also be taken uniformly), the same state is never visited twice within an interval during which the number of variables with status $Z$ is constant. Since the number of variables with status $Z$ is monotonically nonincreasing, the same state is never visited twice overall, and the iterations terminate after finitely many steps.
\end{proof}

That is, for anti-cycling it suffices to apply the classical lexicographic rule only on the ratio-test side, and the lexicographic order $\succ$ with respect to $M$ can be used as is.

\section{Decomposition into the slope problem and the intercept problem}
\label{sec:two-phase}

From Section~\ref{sec:M-infty} onward, in the tests (a$'$)(b$'$) we have compared affine functions $f^1M+f^0$ by the lexicographic order $\succ$ on the coefficient pairs $(f^1,f^0)$. In this section we show that in the main iterations there is essentially no situation in which the two components act simultaneously, and that the iterations split exactly into a ``phase in which comparisons use only the slope $f^1$'' and a ``phase in which the slope stays fixed and comparisons use the intercept $f^0$''.

\subsection{The slope problem}
\label{sec:slope}
For each variable $j$, define the slopes of the bounds by
\[
  l_j^1 := \begin{cases} -1 & (l_j=-\infty)\\ 0 & (l_j>-\infty)\end{cases},\qquad
  u_j^1 := \begin{cases} 1 & (u_j=+\infty)\\ 0 & (u_j<+\infty)\end{cases}
\]
and let $l_j^0$ be $l_j$ if $l_j>-\infty$ and $0$ if $l_j=-\infty$; define $u_j^0$ analogously (for slack variables, $l^1=u^1=l^0=u^0=0$). Then \eqref{eq:trunc} can be written as $\hat l_j(M)=l_j^0+l_j^1M$, $\hat u_j(M)=u_j^0+u_j^1M$. Hence problem~\eqref{eq:lpM}, written with the zero-fixed slack variables $s$ of Section~\ref{sec:init-basis}, is
\begin{equation}
  \min\ c^\top x \quad\text{s.t.}\quad Ax+s=b,\ \ 0\le s\le 0,\ \ l^0+l^1 M\le x\le u^0+u^1 M
  \label{eq:lpM-slack}
\end{equation}
We set the right-hand side $b$ and the constant parts $l^0,u^0$ to $0$, keep only the coefficients $l^1,u^1$ of $M$ as bounds, and consider the problem
\begin{equation}
  \min\ c^\top x^1 \quad\text{s.t.}\quad Ax^1+s=0,\ \ 0\le s\le 0,\ \ l^1\le x^1\le u^1
  \label{eq:slope}
\end{equation}
which we call the \emph{slope problem}. We denote the variables by $x^1$ because, as we will see in Lemma~\ref{lem:basic-slope}, its basic solution coincides exactly with the slope $x^1$ of Lemma~\ref{lem:affine}. In both \eqref{eq:lpM-slack} and \eqref{eq:slope} we keep $s$ as variables, since the combinatorial state $(B,\eta)$ of the algorithm is defined on the indices of $(x,s)$ and the basis $B$ may contain slack variables. Moreover, $0\le s\le0$ are bounds, and in a basic solution that is not primal feasible a basic $s_i$ may take a nonzero value.

The slope problem itself is not new: it coincides with the auxiliary problem that minimizes the sum of dual infeasibilities, known as Phase~1 of the dual simplex method~\cite{fourer1994,kostina2002,maros2003pl,koberstein2005,koberstein2007}. Indeed, the auxiliary problem of Fourer~\cite{fourer1994} is the problem in which the bounds on the unbounded sides are set to $-1,+1$, respectively, and all other bounds and the right-hand side are set to $0$; its bounds are precisely the definition of $l^1,u^1$ (Koberstein~\cite{koberstein2005} removes the variables that are finite on both sides from the auxiliary problem, but these are fixed to $[0,0]$ in the slope problem, so the two are equivalent). In these references, the auxiliary problem is derived directly as the minimization of the sum of dual infeasibilities and is treated as a Phase~1 technique separate from the artificial bounding method~\cite{koberstein2005}. The derivation in this section connects the two by showing that this auxiliary problem arises as the coefficient of $M$ (the slope) extracted from problem~\eqref{eq:lpM} of the artificial bounding method.

The slope problem is the recession cone of the original problem cut by the $\|\cdot\|_\infty$ ball, and it is always feasible at $x^1=0$. Indeed, let $F := \{\,x:\ Ax=b,\ l\le x\le u\,\}$ be the feasible region of the original problem~\eqref{eq:lp} (after preprocessing), and set
\[
  K_j :=
  \begin{cases}
    \{0\} & (\text{finite on both sides})\\
    [0,+\infty) & (\text{only the lower bound finite})\\
    (-\infty,0] & (\text{only the upper bound finite})\\
    \R & (\text{free variable})
  \end{cases}
  ,\qquad
  C \ :=\ \{\,d:\ Ad=0,\ d_j\in K_j\ (j=1,\dots,n)\,\}
\]
$C$ is the cone obtained from $F$ by replacing the right-hand side and the finite bounds by $0$, and when $F\neq\emptyset$ it coincides with the recession cone $\mathrm{rec}(F)$ of $F$. Since $[l_j^1,u_j^1]=K_j\cap[-1,1]$ and, because $s=0$, $Ax^1+s=0$ is equivalent to $Ax^1=0$, the feasible region of \eqref{eq:slope} is $C\cap\{\,d:\|d\|_\infty\le1\,\}$. Writing $x=Mx^1+x^0$ in \eqref{eq:lpM}, dividing by $M$, and letting $M\to\infty$ yields \eqref{eq:slope}; thus \eqref{eq:slope} is the problem that extracts the ``slope'' of \eqref{eq:lpM} as $M\to\infty$. It is well known~\cite{bertsimas1997} that, when $F\neq\emptyset$, ``the original problem is unbounded $\iff$ there exists $d\in\mathrm{rec}(F)$ with $c^\top d<0$'', which is consistent with the distinction $z^1<0$ / $z^1=0$ in Proposition~\ref{prop:trichotomy}(ii)(iii).

\begin{lemma}[Relation between the basic solutions of \eqref{eq:lpM-slack} and \eqref{eq:slope}]
\label{lem:basic-slope}
Fix a combinatorial state $(B,\eta)$, and write the basic solution of \eqref{eq:lpM-slack} for this state as $x(M)=x^0+Mx^1$ (Lemma~\ref{lem:affine}). Then the slope $x^1$ coincides with the basic solution of \eqref{eq:slope} for the same state, and the intercept is $x^0=x(0)$ (both are uniquely determined by $(B,\eta)$).
\end{lemma}

\begin{proof}
Since \eqref{eq:lpM-slack} and \eqref{eq:slope} share the constraint matrix $[\,A\mid I\,]$, the same $(B,\eta)$ determines a basic solution in both. The bounds are $l^0+l^1 M,\,u^0+u^1 M$ and $l^1,u^1$, respectively.

Let $\xi=(\xi_N,\xi_B)$ be the basic solution of \eqref{eq:slope}.

Nonbasic part: if $\eta(j)=L$, then $x_j(M)=l_j^0+l_j^1M$ and $\xi_j=l_j^1$. If $\eta(j)=U$, then $x_j(M)=u_j^0+u_j^1M$ and $\xi_j=u_j^1$. If $\eta(j)=Z$, then $x_j(M)=\xi_j=0$. Therefore $x_N(M)=x_N^0+M\xi_N$, and $\xi_N=x_N^1$. Indeed, $l_j^1=-1$ holds only when $l_j=-\infty$, and $u_j^1=1$ only when $u_j=+\infty$, so $\xi_j$ coincides with the definition of $x_j^1$ in Lemma~\ref{lem:affine}.

Basic part: the right-hand side of the equality constraints of \eqref{eq:lpM-slack} is $b$, and that of the equality constraints $Ax^1+s=0$ of \eqref{eq:slope} is $0$, so
\[
  x_B(M) \ =\ A_B^{-1}\bigl(b-A_Nx_N^0\bigr) + M\,A_B^{-1}\bigl(0-A_N\xi_N\bigr) \ =\ x_B^0 + M \xi_B
\]
and $\xi_B=-A_B^{-1}A_Nx_N^1=x_B^1$. Hence $\xi=x^1$, and setting $M=0$ gives $x^0=x(0)$.
\end{proof}

That is, once $(B,\eta)$ is fixed, the basic solution is a linear function of (right-hand side, bounds); since the data of \eqref{eq:lpM-slack} decompose as $(b,l^0,u^0)+M(0,l^1,u^1)$, the basic solution of problem \eqref{eq:slope}, whose data are the coefficient part $(0,l^1,u^1)$ of $M$, is the slope of $x(M)$. Consequently:
\begin{itemize}
\item The slope $v_i^{\pm,1}$ of the violation $v_i^\pm=v_i^{\pm,1} M+v_i^{\pm,0}$ in \eqref{eq:lpM} equals the violation of row $i$ in \eqref{eq:slope}.
\item The reduced costs and the ratios $\hat c_j/|\hat\alpha_j|$ are determined only by the cost $c$ and the basis $B$, and are therefore common to both problems.
\item The slope $z^1=c^\top x^1$ of $z(M)=z^0+z^1M$ equals the objective value of the basic solution of \eqref{eq:slope}.
\end{itemize}

\subsection{Two-phase decomposition}
\label{sec:two-phase-prop}

\begin{proposition}[Two-phase decomposition]
\label{prop:two-phase}
Consider the combinatorial state transitions of the symbolic artificial bounding method, and distinguish cases according to which of
\[
  \text{Phase A}:\ \exists i,\ v_i^{-,1}>0\ \text{or}\ v_i^{+,1}>0,\qquad
  \text{Phase B}:\ \forall i,\ v_i^{-,1}\le0\ \text{and}\ v_i^{+,1}\le0
\]
the state $(B,\eta)$ is in.
\begin{enumerate}[label=(\roman*)]
\item In Phase A, select the leaving row using only the slopes $w_i^1$ (ties among rows with the largest slope are broken by an arbitrary deterministic rule), and let $w_r=w_r^1M+w_r^0$ be the violation of the selected leaving row $r$ (that is, $w_r^1$ is whichever of $v_r^{-,1},v_r^{+,1}$ is positive; since $v_r^{-,1}+v_r^{+,1}=l_{\beta(r)}^1-u_{\beta(r)}^1\le0$, the two are never positive simultaneously). Then, without using the second component of $\succ$ at all, each iteration is valid as an iteration of the dual simplex method for \eqref{eq:lpM} with sufficiently large $M$. This iteration is exactly the dual simplex method for \eqref{eq:slope}, and $\mathrm{Eligible}=\emptyset$ never occurs. Phase A terminates after finitely many iterations, and $z^1$ at termination equals the optimal value of \eqref{eq:slope}.
\item Once the state enters Phase B, in all subsequent iterations the slopes of the values and bounds of all variables, and hence all $v_i^{\pm,1}$ and $z^1$, remain unchanged, and the state never returns to Phase A.
\end{enumerate}
\end{proposition}

\begin{proof}
(i): The row $r$ selected in Phase A satisfies $w_r^1>0$, so $w_r\succ0$, that is, $r\in V_\infty$. In the dual simplex method, choosing any violated row as the leaving row preserves dual feasibility and does not decrease the dual objective; hence choosing this $r$ in (b$'$) instead of the maximal element with respect to $\succ$ is still valid as an iteration of the dual simplex method for \eqref{eq:lpM} with sufficiently large $M$. The ratio test does not depend on $M$. Since all of the above tests are carried out using only quantities of \eqref{eq:slope}, this iteration is precisely the dual simplex method for \eqref{eq:slope}. Since \eqref{eq:slope} is feasible at $x^1=0$, $\mathrm{Eligible}=\emptyset$, which would certify its infeasibility, does not occur. Anti-cycling follows from the same argument as in Section~\ref{sec:anticycling} (the objective value $z^1$ of \eqref{eq:slope} increases strictly lexicographically with respect to $\epsilon$), and at termination primal feasibility of \eqref{eq:slope}, that is, $v_i^{\pm,1}\le0$ for all $i$, holds.

(ii): The row selected in Phase B satisfies $w_r^1=0$, $w_r^0>0$. The primal step length in the pivot is $w_r/|\hat\alpha_q|$, whose slope is $0$, so neither the slopes $x_B^1$ of the basic variables nor the slope $x_q^1$ of the entering variable $q$ change. The leaving variable $\beta(r)$ is fixed at the bound on the violated side, but since $w_r^1=0$, the slope of that bound equals the slope of the value of $\beta(r)$. Therefore the slopes of the values of all variables are unchanged, and the slopes of the bounds are of course unchanged, so each $v_i^{\pm,1}$ and $z^1=c^\top(\text{slope vector})$ are also unchanged. For the entering variable $q$, its slope $x_q^1$ equals $l_q^1$ or $u_q^1$ ($0$ if in status $Z$), so the slopes $l_q^1-x_q^1$, $x_q^1-u_q^1$ of the new row are both $\le0$. Hence Phase B is preserved.
\end{proof}

\begin{corollary}[Classification obtained at the end of Phase A]
\label{cor:phaseA}
For $z^1$ at the end of Phase A (the optimal value of the slope problem~\eqref{eq:slope}), the following hold.
\begin{enumerate}[label=(\roman*)]
\item If $z^1<0$, then the original problem~\eqref{eq:lp} is unbounded or infeasible, and does not have a finite optimal value.
\item If $z^1=0$, then the original problem~\eqref{eq:lp} has a finite optimal value or is infeasible, and is not unbounded.
\end{enumerate}
\end{corollary}

\begin{proof}
By Proposition~\ref{prop:two-phase}(i), $z^1$ is the optimal value of the slope problem, whose feasible region is $C\cap\{d:\|d\|_\infty\le1\}$ as seen in Section~\ref{sec:slope}. Since $C$ is a cone, $z^1<0$ is equivalent to the existence of $d\in C$ with $c^\top d<0$. If $F=\emptyset$ the claim is trivial, so assume $F\neq\emptyset$. Then $C=\mathrm{rec}(F)$. (i): Taking $d\in\mathrm{rec}(F)$ with $c^\top d<0$ and $x\in F$, we have $x+td\in F$ and $c^\top(x+td)\to-\infty\ (t\to\infty)$, so the problem is unbounded. (ii): $c^\top d\ge0$ for all $d\in\mathrm{rec}(F)$. A linear function on a nonempty polyhedron that is nonnegative on the recession cone is bounded below and attains its minimum, so the problem has a finite optimal value.
\end{proof}

That is, Phase A determines whether the original problem can have a finite optimal value (from the dual side, whether the original dual problem is feasible; Section~\ref{sec:dual-two-phase}). On the other hand, whether the original problem is feasible is not determined by Phase A alone. The original problem may be infeasible even when $z^1=0$, and likewise when $z^1<0$ (the example in Remark~\ref{rem:not-delta}). This distinction is made in Phase B (Proposition~\ref{prop:trichotomy}(i)), and since $z^1$ does not change throughout Phase B (Proposition~\ref{prop:two-phase}(ii)), the classification of Phase A does not contradict the final classification. In particular, if the original problem is known separately to be feasible, unboundedness is determined at the end of Phase A, and Phase B need not be executed.

\subsection{The intercept problem}
\label{sec:intercept}
By Proposition~\ref{prop:two-phase}(ii), the slopes of the values of all variables are constant in Phase B. That is, the slope $x^1=(x_N^1,x_B^1)$ is constant throughout Phase B. At the end of Phase A, $x^1$ is a feasible solution of \eqref{eq:slope}, so $Ax^1=0$ and $l^1\le x^1\le u^1$. Fixing this $x^1$, writing $x=x^0+Mx^1$, and rewriting the constraints of \eqref{eq:lpM} for sufficiently large $M$ as conditions on $x^0$, we obtain the problem
\begin{equation}
  \min\ c^\top x^0 \quad\text{s.t.}\quad Ax^0+s=b,\ \ 0\le s\le 0,\ \ l^B\le x^0\le u^B
  \label{eq:intercept}
\end{equation}
\[
  l^B_j :=
  \begin{cases}
    -\infty & (x_j^1>l_j^1)\\
    l_j & (x_j^1=l_j^1=0)\\
    0 & (x_j^1=l_j^1=-1)
  \end{cases}
  ,\qquad
  u^B_j :=
  \begin{cases}
    +\infty & (x_j^1<u_j^1)\\
    u_j & (x_j^1=u_j^1=0)\\
    0 & (x_j^1=u_j^1=1)
  \end{cases}
\]
which we call the \emph{intercept problem}. Indeed, if $x_j^1>l_j^1$, then $x_j=x_j^0+Mx_j^1$ automatically satisfies the lower bound $\hat l_j(M)=l_j^0+l_j^1M$ for sufficiently large $M$, so the lower bound disappears; if $x_j^1=l_j^1=-1$, then $-M\le x_j^0-M$ leaves $x_j^0\ge0$; and if $x_j^1=l_j^1=0$, then $l_j\le x_j^0$ remains as is (similarly for upper bounds). The objective decomposes as $c^\top x=c^\top x^0+Mz^1$, and $z^1=c^\top x^1$ is constant throughout Phase B. \eqref{eq:intercept} is the problem whose feasible region is the set of points that remain feasible for sufficiently large $M$ when the original problem is shifted by $M$ in the direction $x^1$.

\begin{proposition}[Coincidence of Phase B with the classical dual simplex method]
\label{prop:phaseB}
The iterations of the symbolic artificial bounding method in Phase B ((a$'$)(b$'$)(c)(d) of Step~2) coincide, iteration by iteration, with the iterations of the classical dual simplex method of Section~\ref{sec:dual-simplex} applied to problem~\eqref{eq:intercept}, under the same leaving-row selection rule and tie-breaking. In particular, if $x^1=0$, then \eqref{eq:intercept} is the original problem~\eqref{eq:lp} itself, and Phase B coincides with the classical dual simplex method for the original problem.
\end{proposition}

\begin{proof}
We show that each test of the two iterations sharing the combinatorial state $(B,\eta)$ gives the same result.
\begin{itemize}
\item \textbf{Positions of nonbasic variables}: the slope $x_j^1$ of a nonbasic variable $j$ equals $l_j^1$ or $u_j^1$ ($0$ if in status $Z$), so $j$ lies on a finite bound $x_j^0\in\{l^B_j,u^B_j\}$ of \eqref{eq:intercept}. A variable in status $Z$ satisfies $l_j^1=-1<0=x_j^1<1=u_j^1$, so it is a free variable in \eqref{eq:intercept} as well, and is placed at the value $0$. The basic part is $x_B^0$ (Lemma~\ref{lem:basic-slope}).
\item \textbf{Dual feasibility and the ratio test}: the reduced costs, $\hat c_j$, and the ratios $\hat c_j/|\hat\alpha_j|$ are determined only by the cost $c$ and the basis $B$, and the sign conditions for the statuses $L,U,Z$ are also common, so they coincide in both.
\item \textbf{Identification of violated rows and selection of the leaving row}: in Phase B, $v_i^{\pm,1}\le0$ for all $i$. A side with $v_i^{\pm,1}<0$ is equivalent to the slope of $\beta(i)$ being strictly away from the slope of the bound on that side, that is, to the bound on that side being infinite in \eqref{eq:intercept}, and is not regarded as violated in either. On a side with $v_i^{\pm,1}=0$, the intercept $v_i^{\pm,0}$ itself equals the violation in \eqref{eq:intercept}. Therefore $V_\infty$ coincides with the set of violated rows of \eqref{eq:intercept}, and comparison by $\succ$ reduces to comparison of intercepts, so the selection of the leaving row also coincides.
\item \textbf{Infeasibility test}: $\mathrm{Eligible}=\emptyset$ is the same condition as the infeasibility test of the classical dual simplex method for \eqref{eq:intercept}.
\end{itemize}
Finally, if $x^1=0$, then $x_j^1>l_j^1$ is equivalent to $l_j^1=-1$, that is, $l_j=-\infty$, and $x_j^1=l_j^1=0$ is equivalent to $l_j>-\infty$, so $l^B=l$. Similarly $u^B=u$, and \eqref{eq:intercept} coincides with \eqref{eq:lp}.
\end{proof}

\begin{remark}[Relation to Remark~\ref{rem:absorbing}]
If $x_N^1=0$, then $x_B^1=0$ and hence $x^1=0$, so Remark~\ref{rem:absorbing} corresponds to the case $x^1=0$ of Proposition~\ref{prop:phaseB}. On the other hand, Phase B is entered with $x^1\neq0$ when $z^1<0$, or when $z^1=0$ but nonbasic variables with reduced cost $0$ remain at $\pm M$. In this case \eqref{eq:intercept} is the problem obtained by removing some of the bounds of the original problem and adding bounds $0$ originating from the artificial bounds, and the cleanup of Lemma~\ref{lem:cleanup} can be interpreted as the operation of removing $x^1$ and returning to the original problem. Note that since comparisons in Phase B are only between intercepts, Proposition~\ref{prop:phaseB} also holds for the symbolic artificial bounding method with the original lexicographic comparison.
\end{remark}

\subsection{Overall picture: a chain of problems solved while passing on the basis}
\label{sec:chain}
By Propositions~\ref{prop:two-phase} and~\ref{prop:phaseB}, the symbolic artificial bounding method, and hence the slope--intercept two-phase dual simplex method (Section~\ref{sec:summary}), can be understood as the following chain.
\begin{center}
\begin{tabular}{llll}
\toprule
Stage & Problem solved & Method & Starting basis\\
\midrule
Phase A & Slope problem~\eqref{eq:slope} & Dual simplex method & All-slack basis (dual feasible without computation)\\
Phase B & Intercept problem~\eqref{eq:intercept} & Dual simplex method & Final basis of Phase A\\
Cleanup & Original problem~\eqref{eq:lp} & primal push (Lemma~\ref{lem:cleanup}) & Final basis of Phase B\\
\bottomrule
\end{tabular}
\end{center}
If Phase B stops with infeasibility, the original problem is infeasible; if it stops at an optimum and $z^1<0$, the original problem is unbounded; in either case the method terminates there. When $z^1=0$, if $x^1=0$ the intercept problem is the original problem itself, so no cleanup is needed (Proposition~\ref{prop:phaseB}), and if $x^1\neq0$ the cleanup removes $x^1$. That is, the slope--intercept two-phase dual simplex method is equivalent to the procedure ``solve two LPs in sequence by the dual simplex method and, if necessary, move to a basis of the original problem by crossover''.

Passing on the basis is valid because the three problems share the constraint matrix $A$ and the cost $c$, so the sign conditions on the reduced costs (dual feasibility) carry over unchanged, and because the nonbasic variables of the final basis of Phase A lie on finite bounds of the intercept problem (proof of Proposition~\ref{prop:phaseB}). The cleanup preserves the objective value and primal and dual feasibility (Lemma~\ref{lem:cleanup}).

Note, however, the following two points. First, the intercept problem is not determined by the data alone; it depends on the optimal basic solution $x^1$ obtained in Phase A. If the optimal solution of the slope problem is not unique, the intercept problem differs depending on which optimal solution is reached; thus this is not a procedure that solves fixed LPs in sequence, but a chain in which the solution of the preceding stage determines the problem of the next stage. Second, the equivalence for Phase A is exact for the variant that compares only slopes (Proposition~\ref{prop:two-phase}(i)); in the version with the original lexicographic comparison, discriminating between rows with tied slopes by their intercepts corresponds to a particular choice of tie-breaking in the dual simplex method for the slope problem. Conversely, if ties in the slopes are broken arbitrarily in Phase A, exact agreement with the iteration sequence for a sufficiently large numerical $M$ (Proposition~\ref{prop:consistency}) is in general lost, but dual feasibility, anti-cycling, and the classification into the three cases still hold.

\begin{remark}[Phase A may be replaced by the primal simplex method]
\label{rem:phaseA-primal}
What is passed from Phase A to Phase B in the chain is only the optimal basis of the slope problem~\eqref{eq:slope}, that is, a state $(B,\eta)$ that may include slack variables, which is primal feasible for \eqref{eq:slope} (the Phase B condition $v_i^{\pm,1}\le0$) and dual feasible for \eqref{eq:lpM}. Since Proposition~\ref{prop:two-phase}(ii), Proposition~\ref{prop:phaseB}, and Corollary~\ref{cor:phaseA} use only this property, it does not matter by which method that basis is obtained, and the primal simplex method may be used instead of the dual simplex method. Indeed, under the slack basis $B_0$, if each nonbasic variable is placed on the side with slope $0$ (on the finite side if one exists, and free variables in status $Z$), the basic solution is $x^1=0$, $s=0$, and since the right-hand side of \eqref{eq:slope} is $0$, this is a primal feasible basis. Hence the primal simplex method can start without a Phase~1 and reaches an optimal basis after finitely many iterations. In that basis, the reduced costs of the nonbasic variables with $l_j^1<u_j^1$ satisfy the sign conditions according to the statuses $L,U,Z$. No sign condition is imposed on the variables finite on both sides, which are fixed to $[0,0]$ in \eqref{eq:slope}, but since their slope is $0$, reassigning $\eta(j)$ according to the sign of the reduced cost (Remark~\ref{rem:fixed}; the same operation as (1) in Section~\ref{sec:known-two-phase}) makes the basis dual feasible for \eqref{eq:lpM} without changing $x^1$. In this case, however, the iterations of Phase A are no longer iterations of the dual simplex method for \eqref{eq:lpM}, and the view of the slope--intercept two-phase dual simplex method as a single symbolic dual simplex method for \eqref{eq:lpM} (the symbolic artificial bounding method) holds only for Phase B and the cleanup. Moreover, since the optimal solution of the slope problem is not necessarily unique, the $x^1$ reached, and hence the intercept problem, generally differ from those obtained with the dual simplex method. Note also that at the starting point $x^1=0$ all basic variables are $0$ and primal degeneracy is severe, so which approach is more efficient depends on the problem; we do not evaluate this in this paper.
\end{remark}

\begin{remark}[The separation does not essentially change the variable selection rules]
\label{rem:pricing}
The fact that the problem separates into the slope problem and the intercept problem does not yield new variable selection rules specific to each problem. Both the steepest-edge weights $\gamma_i=\sum_{j\in N}\alpha_{ij}^2$ and the approximate weights of the Devex rule are determined only by the constraint matrix $[\,A\mid I\,]$ and the basis $B$ (and the pivot history), and do not depend on the right-hand side or the bounds. Since the three problems \eqref{eq:lpM-slack}, \eqref{eq:slope}, and \eqref{eq:intercept} share the constraint matrix, the weights at the same basis coincide for the three problems. Therefore ``the steepest-edge rule for the slope problem'' is nothing but the normalization of the slope violations $w_i^1$ by the steepest-edge weights of the original problem, and ``the steepest-edge rule for the intercept problem'' is the normalization of the violations of the intercept problem by the same weights. The same applies to the ratios $\hat c_j/|\hat\alpha_j|$ for the entering variable, since they are also determined only by $c$ and $B$. There is the implementation advantage that the weights can be carried over across phases as they are and no recomputation is required when switching phases, but one cannot expect the separation to improve the quality of selection (there is room to reset the Devex reference framework at the start of Phase B, but its effect would need to be verified empirically).
\end{remark}

\begin{remark}[On the boundary between the phases]
\label{rem:not-delta}
The boundary between Phases A and B in Proposition~\ref{prop:two-phase} differs both from ``while $M$ remains ($x^1\neq0$)'' and from ``while $z^1<0$''.
\begin{enumerate}[label=(\alph*)]
\item Even while $M$ remains, there can be rows with $w_i^1=0$, $w_i^0>0$. For example, in
\[
  \min\ x_1\quad\text{s.t.}\quad x_2=5,\ \ x_1\ \text{free},\ \ 0\le x_2\le1
\]
$x_1=-M$ ($x_1^1=-1$) in the initial state, but the violation of row 1 is $v_1^+=5$ (slope $0$). If (a$'$) is tested using only the first component, then $V_\infty=\emptyset$ and $z^1=-1<0$, and the infeasible problem is misclassified as ``unbounded''.
\item Even if $z^1=0$ occurs in the middle of Phase A, Phase A may continue due to dual degenerate pivots. Conversely, since $z^1$ is fixed in Phase B, if $z^1<0$ at the end of Phase A, Phase B proceeds with $z^1<0$, and whether the problem is infeasible or unbounded is determined in Phase B. The substance of the dual Phase~1 is Phase A, that is, the process of solving \eqref{eq:slope}.
\end{enumerate}
In the cleanup of Lemma~\ref{lem:cleanup}, on the other hand, intercepts are needed when the slope of $t_i$ coincides with that of $M$ or when comparing $t_i$'s with equal slopes, and the lexicographic comparison using both components is indispensable.
\end{remark}

\subsection{Comparison with existing two-phase dual simplex methods}
\label{sec:known-two-phase}

\paragraph{Existing method.}
The slope problem solved in Phase A coincides with the auxiliary problem known as Phase~1 of the dual simplex method~\cite{fourer1994,kostina2002,maros2003pl,koberstein2005,koberstein2007} (Section~\ref{sec:slope}). The existing two-phase dual simplex method using this auxiliary problem (e.g., Algorithm~8 of \cite{koberstein2005}) can be described as follows.
\begin{enumerate}[label=(\arabic*)]
\item Variables finite on both sides are placed at one of their bounds according to the sign of the reduced cost (dual feasibility correction). As a result, the dual infeasibilities handled in Phase~1 are only those arising from variables with an unbounded side.
\item Construct the auxiliary problem in which the right-hand side is set to $0$ and the bounds of variables with an unbounded side are replaced by $[0,1]$ for variables with only a finite lower bound, $[-1,0]$ for variables with only a finite upper bound, and $[-1,1]$ for free variables (variables finite on both sides are removed). Since all variables of this problem are bounded, it can be solved directly, starting from any basis, by an ordinary dual simplex method (Phase~2 code) equipped with the steepest-edge rule and the BFRT. Its optimal value $z_0$, with the sign reversed, is the minimum of the sum of infeasibilities of the original dual problem.
\item If $z_0<0$, the original problem is dual infeasible, and the method stops. In this case it is not distinguished whether the original problem is unbounded or infeasible, and making this distinction requires a separate procedure such as the primal simplex method. For example, when the dual simplex method of HiGHS (version 1.15.1) stops in this state, under the default settings it runs the primal simplex method anew to distinguish the two~\cite{huangfu2018}.
\item If $z_0=0$, the obtained basis is dual feasible for the original problem. Restore the original bounds and right-hand side. The nonbasic variables that were at the bounds $\pm1$ in the auxiliary problem have reduced cost $0$, so they may be placed on the finite side (for free variables, in the nonbasic status with value $0$). Then apply the dual simplex method (Phase~2) to the original problem.
\end{enumerate}
That is, the existing method is a two-phase method consisting of Phase~1 with the auxiliary problem and Phase~2 for the original problem, and $M$ appears nowhere. The artificial bounding method is treated as a Phase~1 technique separate from this~\cite{koberstein2005}.

\paragraph{Differences from Algorithm~\ref{alg:slope-intercept}.}
Phase A solves the same problem as the Phase~1 in (2). The differences between the two appear as follows, depending on $z^1$ and $x^1$ at the end of Phase A.
\begin{enumerate}[label=(\alph*)]
\item If $z^1=0$ and $x^1=0$, the intercept problem is the original problem itself (Proposition~\ref{prop:phaseB}), and the two coincide.
\item If $z^1=0$ and $x^1\neq0$, each term of $z^1=c^\top x^1=\sum_{j\in N}\bar c_jx_j^1$ is nonpositive by dual feasibility, so all nonbasic variables with $x_j^1\neq0$ (variables on artificial bounds) satisfy $\bar c_j=0$. The existing method repositions these on the finite side in (4) to make $x^1=0$, and proceeds to Phase~2 for the original problem. Algorithm~\ref{alg:slope-intercept} leaves them on the artificial bounds, solves the intercept problem, and finally removes $x^1$ by the cleanup of Lemma~\ref{lem:cleanup}. The basis and the basic solution $x_B^0$ at the start of Phase B are the same in both (owing to preprocessing, the intercept values $l_j^0,u_j^0$ of the variables on artificial bounds and the finite bounds they are repositioned to are all $0$), but the problems solved in Phase B differ in their bounds on the support of $x^1$, and the iterations and the final optimal basis generally differ. The optimal value is $z^0$ in both cases.
\item If $z^1<0$, the existing method stops in (3) without distinguishing unboundedness from infeasibility. Algorithm~\ref{alg:slope-intercept} solves the intercept problem by the dual simplex method directly from the basis of Phase A, and makes this distinction using the dual simplex method alone (Proposition~\ref{prop:trichotomy}). If only the distinction were needed, one could instead apply the dual simplex method from the basis of Phase A to the problem with cost $0$ (for which every basis is dual feasible) and test primal feasibility. Algorithm~\ref{alg:slope-intercept} requires no such case distinction on the sign of $z^1$: with the same Phase B procedure as in (a) and (b), it performs classification and solution simultaneously. To the best of our knowledge, no reference describes this treatment.
\end{enumerate}

\paragraph{A unified view.}
The two methods in (b) can be understood in a unified way within the framework of \eqref{eq:lpM}. In \eqref{eq:lpM} with sufficiently large $M$, a nonbasic variable with reduced cost $0$ is dual feasible at either its upper or lower bound, and how it is placed is a degree of freedom inherent in the dual simplex method (an operation of the same kind as a flip in the BFRT or the dual feasibility correction in (1)). If the variables on artificial bounds are left as they are and we proceed to Phase B, then by Proposition~\ref{prop:phaseB} Phase B is the dual simplex method for the intercept problem, and cleanup is required. This is Algorithm~\ref{alg:slope-intercept}. If at the boundary between the phases these are flipped to the opposite finite bound (for free variables, placed in status $Z$), then $x^1=0$, and by Proposition~\ref{prop:no-return} and Remark~\ref{rem:absorbing} Phase B becomes the classical dual simplex method for the original problem, and no cleanup is required. This is the existing method. That is, when $z^1=0$, the two are executions of the same symbolic artificial bounding method that differ only in how the dual degenerate nonbasic variables are placed, and the existing method is derived by adding a single flip to the framework of this paper. The derivation in the reverse direction is asymmetric: from within the existing framework, which has no $M$, neither the intercept problem nor the cleanup arises naturally. For them to arise, a parametrization in $M$ is needed, namely that the optimal solution takes the form $x^0+Mx^1$. Moreover, when $z^1<0$, variables with $\bar c_j\neq0$ remain on the artificial bounds, and flipping them would destroy dual feasibility, so the repositioning cannot be applied.

\paragraph{Scope of novelty.}
From the above, the slope--intercept two-phase dual simplex method solves the same problem as the existing Phase~1 in Phase A, but differs from the existing method as an algorithm in that it solves the intercept problem in Phase B (the contributions of this paper are summarized in Section~\ref{sec:contribution}). In case (b), which of the existing repositioning and the path of this paper is more efficient is not evaluated in this paper, and no theoretical superiority of either is established.

\section{Interpretation from the dual problem}
\label{sec:dual-view}

\subsection{The big-$M$ method for the dual problem}
\label{sec:dual-bigM}

Viewed from the dual side, the construction so far is nothing but the application of the classical big-$M$ method to the dual problem of the original problem, with $M$ treated symbolically. In this section we summarize this correspondence.

Let $J^-:=\{j:l_j=-\infty\}$ and $J^+:=\{j:u_j=+\infty\}$ be the sets of indices with an unbounded side. Restating the bounded problem~\eqref{eq:lpM},
\begin{equation}
  \min\ c^\top x \quad \text{s.t.}\quad Ax=b,\ \hat l\le x\le \hat u
  \tag{\ref{eq:lpM}}
\end{equation}
where, by \eqref{eq:trunc}, $\hat l_j=-M\ (j\in J^-)$, $\hat u_j=M\ (j\in J^+)$, and otherwise $\hat l_j=l_j$, $\hat u_j=u_j$. With dual variables $\kappa,\nu\in\R^n$ for the lower- and upper-bound constraints, the dual problem of this problem is given by
\begin{align}
  \max\quad & b^\top y+\sum_{j\notin J^-} l_j\kappa_j-\sum_{j\notin J^+} u_j\nu_j\ -\ M\Bigl(\sum_{j\in J^-}\kappa_j+\sum_{j\in J^+}\nu_j\Bigr) \label{eq:dualM}\\
  \text{s.t.}\quad & A^\top y+\kappa-\nu=c,\quad \kappa,\nu\ge0 \notag
\end{align}
(the zero-fixed slacks $s$ do not change the primal problem, so they may be omitted here). On the other hand, since the original problem~\eqref{eq:lp} has no bound constraints on the unbounded sides, its dual problem is \eqref{eq:dualM} with
\[
  \kappa_j=0\ \ (j\in J^-),\qquad \nu_j=0\ \ (j\in J^+)
\]
imposed (so that the penalty term vanishes), and it has the dual constraints $\bar c_j\le0\ (j\in J^-)$, $\bar c_j\ge0\ (j\in J^+)$ on $\bar c_j:=c_j-A_j^\top y$ ($\bar c_j=0$ for free variables). That is, $\kappa_j\ (j\in J^-)$ and $\nu_j\ (j\in J^+)$ are \emph{artificial variables} that absorb the violations of these dual constraints, and \eqref{eq:dualM} is exactly the big-$M$ problem obtained by introducing artificial variables into the original dual problem and imposing a penalty with coefficient $M$. In other words, replacing the unbounded bounds of the primal problem by $\mp M$ is the operation dual to relaxing the constraints of the dual problem by a penalty with coefficient $M$. This is the primal--dual counterpart of the well-known fact that the dual of the big-$M$ method applied to a primal problem is the problem in which the dual variables are bounded by $M$.

Furthermore, the dual simplex method for the primal problem is the primal simplex method for the dual problem carried out on the tableau of the primal problem~\cite{lemke1954,bertsimas1997}. Indeed, for a combinatorial state $(B,\eta)$, if we set the basic dual solution determined by complementarity as
\[
  y^\top=c_B^\top A_B^{-1},\qquad
  \kappa_j=\bar c_j\ \ (j\in N,\ \eta(j)=L),\qquad
  \nu_j=-\bar c_j\ \ (j\in N,\ \eta(j)=U)
\]
(with the other components $0$), then the equality constraints $A^\top y+\kappa-\nu=c$ are automatically satisfied, and $\kappa,\nu\ge0$ coincides exactly with the dual feasibility condition of Section~\ref{sec:dual-simplex}. The objective value of \eqref{eq:dualM} is then $b^\top y+\sum_{j\in N}\bar c_jx_j(M)$, which equals $z(M)=z^0+z^1M$ seen in Section~\ref{sec:M-infty}, and its penalty part is
\[
  -z^1 \ =\ \sum_{\substack{j\in N\cap J^-\\ \eta(j)=L}}\kappa_j\ +\sum_{\substack{j\in N\cap J^+\\ \eta(j)=U}}\nu_j ,
\]
that is, the sum of the values of the artificial variables (the infeasibility with respect to the original dual problem). Therefore our algorithm coincides with the big-$M$ method with $M$ treated symbolically, carried out by the primal simplex method on the dual problem~\eqref{eq:dualM}. The components correspond as follows.

\begin{itemize}
\item \textbf{Trivial initial basis} (Section~\ref{sec:init-basis}): at $y=0$ we have $\bar c=c$, the artificial variables are positive exactly for $j\in S$, and their values are $|c_j|$. This is the same structure as in the big-$M$ method, where a trivial initial basis is obtained by letting artificial variables absorb constraint violations. However, artificial variables are needed not for all rows but only for the violated dual constraints $j\in S$. Also, there is no need to add columns for the artificial variables explicitly; placing the nonbasic variables of the primal problem at $\mp M$ represents them implicitly.
\item \textbf{Nonbasic variables on unbounded bounds}: a nonbasic variable $j$ with $x_j^1\neq0$ corresponds to the artificial variable $\kappa_j$ or $\nu_j$ (with value $|\bar c_j|$) remaining in the basis of the dual problem.
\item \textbf{Lexicographic comparison $\succ$} (Section~\ref{sec:M-infty}): the lexicographic comparison of coefficient pairs $(z^1,z^0)$ corresponds to lexicographic optimization~\cite{sherali1983} that gives priority to minimizing the sum of artificial variables (dual infeasibility) over the original objective, and can be regarded as the big-$M$ method treating $M$ as a non-Archimedean infinity~\cite{cococcioni2021} applied to the dual problem. Fixing $M$ as a number gives the artificial bounding method~\cite{koberstein2007,huang2021}. The correspondence between the two components of $\succ$ and the dual Phase~1 and Phase~2 is described in Section~\ref{sec:dual-two-phase}.
\item \textbf{Classification into the three cases} (Proposition~\ref{prop:trichotomy}): $z^1<0$ at termination corresponds to artificial variables remaining positive for all sufficiently large $M$, that is, to the original dual problem being infeasible, and together with $F\neq\emptyset$ this implies unboundedness of the original problem. Termination by $\mathrm{Eligible}=\emptyset$ corresponds to the case where \eqref{eq:dualM} is unbounded. These are the classical tests of the big-$M$ method applied to the dual problem.
\item \textbf{Existence of $M^\dagger$} (Lemma~\ref{lem:union}): the multipliers of the penalized dual constraints are the primal variables $x_j$ themselves, so the argument of Lemma~\ref{lem:union} that $x\in F_M$ if $M\ge\|x\|_\infty$ corresponds to the exact-penalty threshold argument that an $L_1$ penalty becomes exact once the penalty coefficient exceeds the magnitude of the optimal multipliers.
\item \textbf{Cleanup} (Lemma~\ref{lem:cleanup}): an index $j$ with $z^1=0$ and $x_j^1\neq0$ corresponds to an artificial variable remaining in the basis of the dual problem at value $0$, and Lemma~\ref{lem:cleanup} corresponds to the operation of driving it out of the basis. The textbook procedure pivots on any row with a nonzero entry and may therefore lose optimality of the dual problem (feasibility of the primal problem), whereas the primal ratio test of Lemma~\ref{lem:cleanup} selects the row so as to preserve it, and corresponds to the dual push of crossover (Remark~\ref{rem:crossover}).
\item \textbf{No return to unbounded bounds} (Section~\ref{sec:no-return}): reaching $x_N^1=0$ corresponds to all artificial variables having left the basis of the dual problem, after which the iterations are the primal simplex method for the original dual problem, that is, the classical dual simplex method for the original problem (Remark~\ref{rem:absorbing}).
\end{itemize}

\subsection{The two-phase method for the dual problem}
\label{sec:dual-two-phase}
We view the decomposition of Section~\ref{sec:two-phase} from the dual side. We continue to denote the dual variables by $\kappa,\nu$. The dual problem of the slope problem~\eqref{eq:slope} is
\begin{equation}
  \max\ -\Bigl(\sum_{j\in J^-}\kappa_j+\sum_{j\in J^+}\nu_j\Bigr)
  \quad\text{s.t.}\quad A^\top y+\kappa-\nu=c,\ \ \kappa,\nu\ge0
  \label{eq:dual-phase1}
\end{equation}
(since $b=0$ and the finite bounds are $0$, only the dual variables corresponding to unbounded sides appear in the objective). This is the problem obtained from \eqref{eq:dualM} by keeping only the penalty term, that is, exactly the \emph{Phase~1 for the dual problem}, which introduces artificial variables $\kappa_j\ (j\in J^-)$, $\nu_j\ (j\in J^+)$ into the original dual problem and minimizes their sum. Its optimal value equals $z^1$.

On the other hand, the dual problem of the intercept problem is
\begin{align*}
  \max\quad & b^\top y+\sum_{j} l^B_j\kappa_j-\sum_{j} u^B_j\nu_j\\
  \text{s.t.}\quad & A^\top y+\kappa-\nu=c,\ \ \kappa,\nu\ge0,\ \ \kappa_j=0\ (x_j^1>l_j^1),\ \ \nu_j=0\ (x_j^1<u_j^1)
\end{align*}
Since no multiplier corresponds to a bound that does not exist in the primal problem, the dual variables on the sides where the bound is infinite are fixed to $0$. That is, removing bounds in the intercept problem corresponds to fixing the corresponding multipliers to $0$ in the dual problem.

Meanwhile, $x^1$ at the end of Phase A is an optimal solution of the slope problem~\eqref{eq:slope}, that is, of the primal problem of \eqref{eq:dual-phase1}. By the complementary slackness theorem, once an optimal solution of the primal problem is fixed, a feasible solution of the dual problem is optimal if and only if it is complementary to that fixed optimal solution. Therefore a feasible solution $(y,\kappa,\nu)$ of \eqref{eq:dual-phase1} is optimal if and only if it satisfies complementarity with $x^1$, $\kappa_j(x_j^1-l_j^1)=0$, $\nu_j(u_j^1-x_j^1)=0$, that is,
\[
  \kappa_j=0\ \ (x_j^1>l_j^1),\qquad \nu_j=0\ \ (x_j^1<u_j^1)
\]
(no condition is imposed for the slack variables $s$, since $l^1=u^1=x^1=0$ for them). This coincides index by index with the conditions for the multipliers fixed to $0$ in the dual of the intercept problem, and the remaining constraints $A^\top y+\kappa-\nu=c$, $\kappa,\nu\ge0$ are also common. Therefore the feasible region of the dual of the intercept problem coincides exactly with the \emph{optimal face of the dual Phase~1}. Moreover, the artificial variables corresponding to the bounds originating from artificial bounds ($l^B_j=0$ if $x_j^1=l_j^1=-1$, $u^B_j=0$ if $x_j^1=u_j^1=1$) remain on this face as variables with cost $0$ (on the optimal face, the sum of the artificial variables is fixed at $-z^1$). Hence solving the intercept problem corresponds to maximizing the original dual objective over the optimal face of the dual Phase~1, that is, to the \emph{Phase~2 of the dual problem}. The disappearance of some of the bounds of the original problem in the intercept problem is the dual counterpart of the operation, in the two-phase method, of removing the variables with positive reduced cost after Phase~1 (restricting to the optimal face).

Thus the symbolic artificial bounding method is at the same time the big-$M$ method for the dual problem and the two-phase method for the dual problem, and this correspondence holds at the level of iterations, not at the level of limiting objective values. The classical fact that the big-$M$ method coincides with the two-phase method as $M\to+\infty$~\cite{chvatal1983,bertsimas1997} is usually stated as a comparison of objective values, but in our setting the first and second components of $\succ$ are assigned exactly to Phase A (Phase~1) and Phase B (Phase~2), respectively. The algorithm of Section~\ref{sec:summary}, based on this correspondence, executes the two phases in sequence without using the lexicographic comparison.

\section{The slope--intercept two-phase dual simplex method: summary}
\label{sec:summary}

Integrating all of the above components, we summarize the \emph{slope--intercept two-phase dual simplex method} for the general linear optimization problem~\eqref{eq:lp} as Algorithm~\ref{alg:slope-intercept}. The word ``dual'' in the name refers both to the fact that each phase is a dual simplex method and to the fact that the whole corresponds to a two-phase method for the dual problem (Section~\ref{sec:dual-two-phase}). The tests (a$'$)(b$'$) of Section~\ref{sec:M-infty} defined the iterations by comparing affine functions with the lexicographic order $\succ$ on coefficient pairs, but, as seen in Section~\ref{sec:two-phase}, the first component (slope) and the second component (intercept) of this comparison do not act simultaneously in the main iterations. Accordingly, in Algorithm~\ref{alg:slope-intercept} the main iterations are executed not via the lexicographic comparison but as two phases that solve the slope problem~\eqref{eq:slope} and the intercept problem~\eqref{eq:intercept} in sequence. The inputs $A,b,c,l,u$ may be the original data before preprocessing, and free variables require no special preprocessing.

After preprocessing (Section~\ref{sec:shift-preprocess}), Algorithm~\ref{alg:slope-intercept} consists of the following two phases and a cleanup performed when necessary.
\begin{itemize}
\item \textbf{Phase A (slope problem)}: starting from the dual feasible initial basis obtained without computation (Section~\ref{sec:init-basis}), run the dual simplex method on the slope problem~\eqref{eq:slope}. Only the basic solution $x^1$ of the slope problem is maintained; the intercept $x^0$ is not maintained.
\item \textbf{Phase B (intercept problem)}: from the final basis of Phase A, run the ordinary real-valued dual simplex method on the intercept problem~\eqref{eq:intercept}. At termination, determine whether the original problem is infeasible, unbounded, or has a finite optimal value.
\item \textbf{Cleanup}: only when there is a finite optimal value and nonbasic variables remain on artificial bounds ($x_N^1\neq0$), construct an optimal basic solution free of $M$ by the primal push of Lemma~\ref{lem:cleanup}.
\end{itemize}
Finally, the inverse shift returns to the original coordinate system. Ties in the ratio test are broken by the lexicographic rule based on cost perturbation in Section~\ref{sec:anticycling}. Since the problems solved in Phases A and B are both ordinary LPs, in an implementation the standard techniques of the dual simplex method, such as leaving-row selection by the steepest-edge or Devex rules and the Bound Flipping Ratio Test, can be applied directly to each phase without any special treatment of $M$.

\begin{center}
\refstepcounter{algorithm}
\textbf{Algorithm~\thealgorithm.} \textsc{SlopeInterceptDualTwoPhase}$(A,b,c,l,u)$ --- the slope--intercept two-phase dual simplex method
\label{alg:slope-intercept}
\end{center}
\begin{algorithmic}[1]

\Statex \textbf{Preprocessing}
\State Compute the shift $t_j$ for each $j$, and replace $A,b,l,u \gets A,\ b-At,\ l-t,\ u-t$ ($t_j=0$ for free variables) \Comment{Section~\ref{sec:shift-preprocess}}
\State Determine the slopes $l_j^1,u_j^1$ of the bounds for each $j$ ($l^1=u^1=0$ for slack variables) \Comment{Section~\ref{sec:two-phase}}

\Statex \textbf{Phase A: slope problem~\eqref{eq:slope}}
\State $B\gets\{n+1,\dots,n+m\}$ (slack indices); for each $j=1,\dots,n$, following \eqref{eq:init-status}, set $\eta(j)\gets L$ ($c_j>0$, or $c_j=0$ and $l_j>-\infty$), $Z$ ($c_j=0$ and free variable), $U$ (otherwise) \Comment{dual feasible without computation, Proposition~\ref{prop:dual-feasible-init}}
\State $x_N^1\gets$ the nonbasic slopes of Lemma~\ref{lem:affine}; $x_B^1\gets -A_Nx_N^1$ \Comment{$A_B=I$. Basic solution $x^1$ of the slope problem, Lemma~\ref{lem:basic-slope}}
\Loop
  \State For each row $i$, $v_i^{-,1}\gets l_{\beta(i)}^1-x_{B,i}^1$, $v_i^{+,1}\gets x_{B,i}^1-u_{\beta(i)}^1$ \Comment{slopes of the violations}
  \State $V_A\gets\{\,i : v_i^{-,1}>0 \text{ or } v_i^{+,1}>0\,\}$
  \If{$V_A=\emptyset$}
    \State \textbf{break} \Comment{optimal basis of the slope problem reached; go to Phase B}
  \EndIf
  \State With $w_i^1\gets\max(v_i^{-,1},v_i^{+,1})$, $r\gets\operatorname{arg\,max}_{i\in V_A} w_i^1/\sqrt{\gamma_i}$ \Comment{Dantzig's rule: $\gamma_i\equiv1$; ties by an arbitrary deterministic rule}
  \State $\rho\gets+1$ (if $v_r^{-,1}>0$), $-1$ (if $v_r^{+,1}>0$); $\alpha^\top\gets e_r^\top A_B^{-1}A_N$; $\sigma_j,\hat\alpha_j,\hat c_j\gets\cdots$ \Comment{Step~2(c)}
  \State $\mathrm{Eligible}\gets\{j\in N : l_j<u_j,\ \rho\,\hat\alpha_j<0\}$ \Comment{variables finite on both sides are also candidates; never empty}
  \State $q\gets\operatorname{arg\,min}_{j\in\mathrm{Eligible}}\hat c_j/|\hat\alpha_j|$ \Comment{status-$Z$ candidates have ratio $0$ and take top priority; ties by the rule of Section~\ref{sec:anticycling}}
  \State Pivot exchanging $\beta(r)\leftrightarrow q$; update $B,A_B^{-1},x^1,\eta$ \Comment{$\beta(r)$ goes to the bound on the violated side}
\EndLoop

\Statex \textbf{Phase B: intercept problem~\eqref{eq:intercept}}
\State $z^1\gets c^\top x^1$; determine $l^B,u^B$ by \eqref{eq:intercept} \Comment{$x^1$ and $z^1$ are fixed in Phase B}
\State $x_N^0\gets$ ($l^B_j,u^B_j,0$ according to status $L,U,Z$); $x_B^0\gets A_B^{-1}(b-A_Nx_N^0)$ \Comment{one FTRAN}
\Loop \Comment{dual simplex method of Section~\ref{sec:dual-simplex} for \eqref{eq:intercept}}
  \State $V_B\gets\{\,i : x_{B,i}^0<l^B_{\beta(i)} \text{ or } x_{B,i}^0>u^B_{\beta(i)}\,\}$
  \If{$V_B=\emptyset$}
    \State \textbf{break} \Comment{optimal basis of the intercept problem reached}
  \EndIf
  \State $r\gets\operatorname{arg\,max}_{i\in V_B}\Delta_i/\sqrt{\gamma_i}$ ($\Delta_i$ is the violation with respect to $l^B,u^B$); $\rho\gets$ direction of the violation; $\alpha^\top,\sigma_j,\hat\alpha_j,\hat c_j\gets\cdots$
  \State $\mathrm{Eligible}\gets\{j\in N : l^B_j<u^B_j,\ \rho\,\hat\alpha_j<0\}$
  \If{$\mathrm{Eligible}=\emptyset$}
    \State \Return ``the original problem is infeasible'' \Comment{Proposition~\ref{prop:alg}(i)}
  \EndIf
  \State $q\gets\operatorname{arg\,min}_{j\in\mathrm{Eligible}}\hat c_j/|\hat\alpha_j|$; pivot exchanging $\beta(r)\leftrightarrow q$; update $B,A_B^{-1},x^0,\eta$
\EndLoop
\If{$z^1<0$}
  \State \Return ``the original problem is feasible and unbounded'' \Comment{Proposition~\ref{prop:alg}(ii)}
\EndIf
\State $z^0\gets c^\top x^0$ \Comment{Proposition~\ref{prop:alg}(iii), optimal value}

\Statex \textbf{Cleanup}
\State Let the basic solution be $x(M)=x^0+Mx^1$ ($x^1$ from the end of Phase A, $x^0$ from the end of Phase B) \Comment{Lemma~\ref{lem:basic-slope}}
\While{there exists a nonbasic variable $j$ with $x_j^1\neq0$} \Comment{at most $K$ times}
  \State Choose one such $j$, and compute $t^\ast$ by the primal ratio test of Lemma~\ref{lem:cleanup} ($\succ$ comparison of affine functions of $M$)
  \If{$t^\ast=M$}
    \State Place $x_j$ on the finite side (for a free variable, status $Z$, value $0$) \Comment{basis unchanged, Remark~\ref{rem:state-F}}
  \Else
    \State Pivot exchanging $\beta(r)\leftrightarrow j$ at the row $r$ with $t^\ast=t_r$; fix $\beta(r)$ at the finite bound reached
  \EndIf
  \State Update $x^0,x^1,\eta$ (and also $B,A_B^{-1}$ if a pivot was performed) \Comment{objective value, primal feasibility, and dual feasibility all unchanged}
\EndWhile
\State $x^{\ast\ast}\gets$ the resulting basic solution $x^0$ (free of $M$ since $x^1=0$) \Comment{optimal basic solution of the original problem (in the preprocessed coordinates)}
\State \Return $\bigl(x^{\ast\ast}+t,\ \ z^0+c^\top t\bigr)$ \Comment{inverse shift; optimal solution and optimal value in the original coordinates}
\end{algorithmic}

\paragraph{Correctness.}
The correctness of Algorithm~\ref{alg:slope-intercept} can be shown directly, without relying on the theory of lexicographic comparison of Section~\ref{sec:M-infty}. Let $F$ be the feasible region of the original problem~\eqref{eq:lp} (after preprocessing).

\begin{proposition}[Correctness of Algorithm~\ref{alg:slope-intercept}]
\label{prop:alg}
Algorithm~\ref{alg:slope-intercept} terminates after finitely many iterations, and the following hold.
\begin{enumerate}[label=(\roman*)]
\item It returns ``infeasible'' if and only if $F=\emptyset$.
\item It returns ``feasible and unbounded'' if and only if $F\neq\emptyset$ and $\inf_{x\in F}c^\top x=-\infty$.
\item Otherwise, the returned $x^{\ast\ast}+t$ is an optimal basic solution of the original problem, and $z^0+c^\top t$ is its optimal value.
\end{enumerate}
\end{proposition}

\begin{proof}
\emph{Phase A.} The slope problem is an LP in which all bounds are finite, and the initial basis is dual feasible (the correspondence between the signs of the reduced costs $c$ and the statuses is the same as in Proposition~\ref{prop:dual-feasible-init}). Phase A is the dual simplex method for the slope problem. Variables finite on both sides are fixed to $[0,0]$ in the slope problem, but including them as candidates in the ratio test is still valid as a dual simplex method, and preserves the dual feasibility of those variables for Phase B and beyond. Under the lexicographic rule based on cost perturbation, Phase A terminates after finitely many iterations, and since the slope problem is feasible at $x^1=0$, it reaches an optimal basis at termination. $z^1=c^\top x^1$ is its optimal value.

\emph{Phase B.} The reduced costs are determined only by $c$ and the basis $B$, and the sign conditions for the statuses $L,U,Z$ do not depend on the problem, so the final basis of Phase A is dual feasible for the intercept problem as well. Moreover, each nonbasic variable $j$ satisfies $x_j^1=l_j^1$ (status $L$), $x_j^1=u_j^1$ (status $U$), or $x_j^1=0$ (status $Z$, in which case $j$ is a free variable), so by the definition \eqref{eq:intercept}, $j$ lies on a finite bound $l^B_j$ or $u^B_j$ of the intercept problem, and a variable in status $Z$ is a free variable in the intercept problem as well. Therefore Phase B can be executed as an ordinary dual simplex method for the intercept problem, and it terminates after finitely many iterations, either with ``infeasible'' or at an optimal basis.

\emph{Relation between the intercept problem and the original problem.} At the end of Phase A, $x^1$ is a feasible solution of the slope problem, so $Ax^1=0$ and $l^1\le x^1\le u^1$. (a) If $x^0$ is a feasible solution of the intercept problem, then $x:=x^0+Mx^1\in F$ for sufficiently large $M$. Indeed, $Ax=b$; in components with $x_j^1>l_j^1$, $x_j$ satisfies the lower bound ($l_j$ if finite, no constraint otherwise) as $M\to\infty$; in components with $x_j^1=l_j^1=-1$, $x_j^0\ge0$ only gives $x_j\ge-M$, but the lower bound of the original problem is $-\infty$; and in components with $x_j^1=l_j^1=0$, $x_j=x_j^0\ge l_j$ (similarly for upper bounds). (b) Conversely, if $x\in F$, then $x^0:=x-Mx^1$ is a feasible solution of the intercept problem for sufficiently large $M$. Indeed, $Ax^0=b$; in components with $x_j^1=l_j^1=-1$, $x_j^0=x_j+M\ge0$; in components with $x_j^1=l_j^1=0$, $x_j^0=x_j\ge l_j$; and components with $x_j^1>l_j^1$ have no lower bound (similarly for upper bounds). In either case $c^\top x=c^\top x^0+Mz^1$.

(i): By (a)(b), infeasibility of the intercept problem is equivalent to $F=\emptyset$.
(ii)(iii): When Phase B reaches an optimal basis, $F\neq\emptyset$. If $z^1<0$, the original problem is unbounded by Corollary~\ref{cor:phaseA}(i). If $z^1=0$, then by (b), $c^\top x=c^\top(x-Mx^1)\ge z^0$ for any $x\in F$, and by (a), $x^0+Mx^1\in F$ attains this value $z^0$, so $z^0$ is the optimal value of the original problem. In this case the basic solution $x(M)=x^0+Mx^1$ satisfies the bounds of the original problem for all sufficiently large $M$, and the state is dual feasible for the original problem, so the assumptions of Lemma~\ref{lem:cleanup} are satisfied, and the cleanup yields an optimal basic solution $x^{\ast\ast}$ of the original problem while preserving the objective value $z^0$. The final inverse shift is justified by Section~\ref{sec:shift-preprocess}.
\end{proof}

On the other hand, by Section~\ref{sec:two-phase}, each iteration of Phases A and B is also a valid iteration of the dual simplex method for \eqref{eq:lpM} with sufficiently large $M$ (since ties in the slopes are broken arbitrarily in Phase A, exact agreement with the iteration sequence for a numerical $M$ is in general lost; Section~\ref{sec:chain}). That is, Algorithm~\ref{alg:slope-intercept} executes the iterations of the artificial bounding method in the limit $M\to+\infty$, split into the two phases of the slope problem and the intercept problem. Note that when $z^1=0$, the existing repositioning~\cite{fourer1994,koberstein2005} may also be used instead of the cleanup (Section~\ref{sec:known-two-phase}).

The main iterations of Phases A and B do not use the lexicographic comparison $\succ$ of affine functions at all. Phase A compares only slopes (real values with respect to bounds with integer coefficients $0,\pm1$), and Phase B compares only intercepts (ordinary real values). $\succ$ remains only in the primal ratio test of the cleanup, where both components are essentially needed, e.g., when the slope of $t_i$ coincides with that of $M$ (Remark~\ref{rem:not-delta}).

\begin{remark}[Comparison with execution by lexicographic comparison]
\label{rem:impl}
We compare Algorithm~\ref{alg:slope-intercept} with executing the main iterations as a single sequence of iterations using the lexicographic comparison of (a$'$)(b$'$). The main computations per iteration (BTRAN, FTRAN, computation of the pivot row $\alpha^\top$) are unchanged, and the number of iterations is the same except for differences in tie-breaking, so the reduction in computational cost is small. On the other hand, there are the following advantages (their effect on computing time has not been measured separately).
\begin{enumerate}[label=(\alph*)]
\item \textbf{Early switch to real arithmetic}: Phase B is an ordinary real-valued dual simplex method that treats the bounds on sides with $v_i^{\pm,1}<0$ as nonexistent, tests sides with $v_i^{\pm,1}=0$ by constant violations, and does not maintain pairs of affine functions. The switching point of Remark~\ref{rem:absorbing} (reaching $x_N^1=0$) is moved forward to the end of Phase A, and this difference is large for problems that terminate with $x_N^1\neq0$ (unbounded cases and cases requiring cleanup). In Phase A the intercept $x^0$ need not be updated either; it suffices to compute it with one FTRAN at the transition to Phase B.
\item \textbf{Numerical robustness}: the floating-point equality test ``are the slopes equal'' between different rows, which was needed in the lexicographic comparison $\succ$, becomes unnecessary; what remains is only the sign test of $w_i^1$ for each row (of the same kind as the usual primal feasibility tolerance). Moreover, \eqref{eq:slope} has right-hand side $0$ and bounds only $0,\pm1$, and is not affected by the scale of $b$.
\end{enumerate}
As a concern, the effect on the number of iterations of losing the discrimination of tied slopes by intercepts in Phase A needs to be checked empirically.
\end{remark}

\section{Numerical experiments}
\label{sec:experiments}

We compare ENOMOTO-Solver (implemented in Rust, with a Python interface), a solver implementing Algorithm~\ref{alg:slope-intercept}, with the open-source LP solvers HiGHS\cite{huangfu2018}, CLP\cite{clp}, and SoPlex\cite{soplex}. The test set consists of 207 problems: problems with a finite optimal solution from Netlib, Kennington, and Mittelmann, together with infeasible problems and unbounded problems constructed as their duals. The aim of this section is to show that the slope--intercept two-phase dual simplex method runs at a practical speed when combined with standard techniques such as the steepest-edge rule, BFRT, and the Forrest--Tomlin update, and that distinguishing infeasibility from unboundedness takes little additional time.

\subsection{Overview of the implementation}
\label{sec:exp-impl}

ENOMOTO-Solver proceeds as follows.
\begin{enumerate}[label=(\arabic*)]
\item \textbf{Presolve}: scaling by Ruiz equilibration\cite{ruiz2001}, removal of duplicate and dominated rows, bound propagation, row and column singletons, equality doubletons, dual fixing, substitution of implied free columns, merging parallel columns, and so on, repeated until a fixed point is reached.
\item \textbf{Conversion to standard form}: we apply the shift of Section~\ref{sec:shift-preprocess} and attach a zero-fixed slack to each row. Infinite bounds are kept infinite rather than replaced by a numerical $M$. When the constraint matrix splits into several connected components, each component is solved separately.
\item \textbf{Phase A, Phase B, and cleanup}: executed as in Algorithm~\ref{alg:slope-intercept}. If no variable is placed on an infinite side in the initial configuration ($S=\emptyset$), Phase A is skipped. The leaving row is chosen by the dual steepest-edge rule\cite{forrestgoldfarb1992}, the ratio test is BFRT, and the basis matrix is factorized by sparse LU factorization with the Markowitz rule\cite{markowitz1957} and the Forrest--Tomlin update\cite{forrest1972}.
\item \textbf{Final polishing and postsolve}: the resulting basis is checked against the true bounds and costs. Only if dual infeasible columns remain because of floating-point errors do we finish with the primal simplex method. Finally, the presolve is traversed in reverse to recover a solution of the original problem.
\end{enumerate}
In the default setting, if $z^1<0$ at the end of Phase A, the solver returns ``infeasible or unbounded'' and stops (the same treatment as in the existing method of Section~\ref{sec:known-two-phase}). If the solver is configured to proceed to Phase B, it distinguishes the two cases as in Proposition~\ref{prop:alg}.

\subsection{Implementation techniques not found in HiGHS}
\label{sec:exp-features}

By comparison with the source code of HiGHS 1.15.1, we list the main techniques that are absent from, or treated differently in, the HiGHS dual simplex method.\begin{enumerate}[label=(\alph*)]
\item \textbf{Slope--intercept phase structure}: the dual phase 1 of HiGHS also solves the same auxiliary problem as the slope problem, but differs in that it bounds free variables by $[-1000,1000]$. Moreover, after phase 1, HiGHS restores the original bounds and performs phase 2, and if the optimal value of phase 1 is negative, it leaves the dual simplex method and runs the primal simplex method to distinguish infeasibility from unboundedness (Section~\ref{sec:known-two-phase}). Our implementation makes this distinction within the dual simplex method as well, in Phase B; it solves the intercept problem in Phase B and removes variables on artificial bounds in the cleanup.
\item \textbf{Fused FTRAN}: the three FTRANs needed in one iteration (the entering column, the vector used to update the dual steepest-edge weights, and the contribution of the variables flipped by BFRT) are solved together in a single pass over the LU factors. HiGHS solves them separately.
\item \textbf{Bordered LU factorization}: nearly dense columns are separated and handled separately as a Schur complement.
\item \textbf{Reuse of the pivot order}: at refactorization, the column pivot order of the previous factorization is reused as the initial order. In HiGHS, a similar mechanism is used only for hot starts.
\item \textbf{Refactorization criteria}: every few iterations, the residual $\|A_Bx_B-\text{right-hand side}\|$ of the basic solution is checked, and a refactorization is performed when accumulation of errors is detected. In addition, we ported the HiGHS rule that schedules refactorization on the basis of an estimate of the computational cost, and tuned its thresholds on Netlib.
\end{enumerate}
On the other hand, HiGHS also has almost all of the presolve techniques. Parallelization, such as solving connected components in parallel, is also implemented, but its effect is not measured separately in this section (the numbers of threads used are given in Section~\ref{sec:exp-setup}).

\subsection{Experimental environment and measurement}
\label{sec:exp-setup}

The computer is a laptop with an AMD Ryzen 7 5700U (8 cores, 16 threads) and 16\,GB of memory; all measurements were taken under WSL2 (Ubuntu 22.04, memory limit 12\,GB, no swap) on Windows 11. We compare with HiGHS 1.15.1 (Python binding highspy), CLP 1.17.11, and SoPlex 8.1.0. CLP and SoPlex were built from source and called from small programs that use their libraries directly. All solvers were run with their default settings except for the time limit and the number of threads (presolve is on in all of them; CLP chooses the algorithm automatically as by default, and the dual simplex method of HiGHS is the default serial version). ENOMOTO-Solver and HiGHS were given 16 threads, the number of logical CPUs; CLP and SoPlex are serial solvers. ENOMOTO-Solver was run in the mode that distinguishes infeasibility from unboundedness (that is, it proceeds to Phase B).

The test set consists of 207 problems in the following five sets.
\begin{itemize}
\item The 93 LP problems of Netlib\cite{gay1985} that have a finite optimal solution.
\item The 29 infeasible LP problems of Netlib.
\item 29 LP problems obtained from the above 29 problems by taking the dual of the problem with the objective set to $0$. By Farkas' lemma, if the original problem is infeasible, its dual is feasible and unbounded. However, cplex2 is extremely close to being feasible, and its dual is not numerically unbounded: HiGHS, CLP, and SoPlex all report the optimal value $0$. We therefore do not fix an expected answer for this one problem.
\item The 16 problems of the Kennington test set\cite{carolan1990}.
\item 40 problems of Mittelmann's LP benchmarks\cite{mittelmann} (LPopt): of the 44 publicly available problems we could obtain, we excluded the four largest ones that do not fit in the memory of this computer (thk\_48, L2CTA3D, dlr2, and Dual2\_5000, with $2.8\times10^7$--$9.3\times10^7$ nonzeros).
\end{itemize}

We measure the wall-clock time of the solve call; it excludes reading the MPS file and building the model, and includes presolve and postsolve. Every solve is run in a fresh process, and each problem is solved three times by each solver and the median is taken (only once if the first run took 300\,s or more or failed). The order of the solvers is rotated between repetitions, and if the first run took less than 1\,s, it is discarded and measured again. The time limit is 600\,s. A problem counts as solved if the solver returns the expected status and, for an optimal status, the relative difference $|z-z^\ast|/\max(1,|z^\ast|)$ from the reference value is at most $10^{-6}$. As the reference value $z^\ast$ we take the value, among those of the solvers that report optimality, that agrees with the largest number of solvers. We report the geometric mean, the shifted geometric mean with a shift of $10$\,s, $\exp\bigl(\frac1N\sum_i\log(t_i+10)\bigr)-10$, and the total time; unsolved problems are counted at the time limit of 600\,s. As an indication of measurement noise, among the 76 cases measured three times with a median of at least 1\,s, the median of (maximum$-$minimum)/median was 2.2\% and its 90th percentile was 6.6\%.

Table~\ref{tab:summary} summarizes the results for each set.

\begin{table}[htbp]
\centering
\caption{Summary per test set (times in seconds). Unsolved problems are counted at the time limit of 600\,s in the geometric mean, the shifted geometric mean, and the total. Bold marks the best value in each set.}
\label{tab:summary}
\small
\begin{tabular}{llrrrr}
\toprule
Test set & Solver & Solved & Geo.\ mean & Shifted geo.\ mean & Total\\
\midrule
Netlib (93) & ENOMOTO & 93 & 0.0136 & \textbf{0.152} & \textbf{16.0}\\
 & HiGHS & 93 & 0.0297 & 0.166 & 17.3\\
 & CLP & 93 & \textbf{0.0108} & 0.164 & 17.5\\
 & SoPlex & 92 & 0.0230 & 0.741 & 632\\
\midrule
Kennington (16) & ENOMOTO & 16 & 0.303 & \textbf{0.964} & \textbf{17.5}\\
 & HiGHS & 16 & 0.580 & 1.68 & 34.9\\
 & CLP & 16 & \textbf{0.293} & 1.16 & 22.4\\
 & SoPlex & 16 & 1.44 & 5.76 & 208\\
\midrule
Mittelmann (40) & ENOMOTO & 12 & \textbf{367} & \textbf{377} & 19{,}024\\
 & HiGHS & 11 & 401 & 405 & 19{,}200\\
 & CLP & \textbf{18} & 386 & 393 & \textbf{18{,}873}\\
 & SoPlex & 7 & 501 & 503 & 21{,}722\\
\midrule
Infeasible (29) & ENOMOTO & 29 & \textbf{0.00156} & \textbf{0.00597} & \textbf{0.173}\\
 & HiGHS & 29 & 0.00749 & 0.0154 & 0.448\\
 & CLP & 29 & 0.00490 & 0.0743 & 2.26\\
 & SoPlex & 28 & 0.00469 & 1.53 & 600\\
\midrule
Unbounded (29) & ENOMOTO & 29 & 0.00740 & 0.0588 & 1.74\\
 & HiGHS & 29 & 0.0197 & 0.335 & 12.1\\
 & CLP & 28 & 0.0161 & 2.01 & 615\\
 & SoPlex & 29 & \textbf{0.00639} & \textbf{0.0332} & \textbf{0.975}\\
\midrule
All (207) & ENOMOTO & 179 & \textbf{0.0843} & \textbf{10.6} & \textbf{19{,}060}\\
 & HiGHS & 178 & 0.183 & 11.0 & 19{,}265\\
 & CLP & \textbf{184} & 0.100 & 11.3 & 19{,}530\\
 & SoPlex & 172 & 0.146 & 13.4 & 23{,}163\\
\bottomrule
\end{tabular}
\end{table}

ENOMOTO-Solver returned no wrong answer on any of the 207 problems, and on the problems for which it reported optimality, the relative difference from the reference value was at most $9.8\times10^{-11}$. Of the 28 problems it did not solve, 26 exceeded the time limit, one ran out of memory (thk\_63), and on one it gave up for numerical reasons (set-cover-model). The other solvers returned an unexpected answer in three cases: SoPlex returned an optimal value different from the reference on pilot87 and ``optimal'' on the infeasible cplex2, and CLP returned ``optimal'' on the unbounded dual of qual. On the problems solved by both solvers, the geometric mean of the ratio of the time of ENOMOTO-Solver to that of HiGHS was 0.41 (faster on 155 of 177 problems); to CLP, 0.84 (82 of 174); and to SoPlex, 0.58 (133 of 170).

\subsection{Problems with a finite optimal solution}
\label{sec:exp-optimal}

On the 93 Netlib problems, the geometric mean of the ratio of ENOMOTO-Solver to HiGHS was 0.46, and ENOMOTO-Solver was faster on 84 problems. Against CLP, however, the geometric mean of the ratio was 1.26 and ENOMOTO-Solver was faster on only 24 problems. Most Netlib problems are solved within a few milliseconds, and on such problems the small fixed cost of CLP pays off. In the shifted geometric mean and the total time, which weigh the longer problems more heavily, ENOMOTO-Solver is the smallest of the four solvers. Table~\ref{tab:netlib-large} lists the 25 problems on which some solver needed at least 0.1\,s.

\begin{table}[htbp]
\centering
\caption{Netlib problems on which some solver needed at least 0.1\,s (times in seconds, median of three runs). $^\dagger$ indicates an optimal value different from the reference value.}
\label{tab:netlib-large}
\small
\begin{tabular}{lrrrr}
\toprule
Problem & ENOMOTO & HiGHS & CLP & SoPlex\\
\midrule
dfl001 & 5.87 & 5.66 & 4.63 & 13.4\\
pilot87 & 4.02 & 4.18 & 5.98 & ---$^\dagger$\\
fit2p & 0.976 & 0.996 & 1.29 & 2.14\\
pilot & 0.800 & 0.962 & 0.889 & 2.65\\
maros-r7 & 0.673 & 0.794 & 1.15 & 1.83\\
d2q06c & 0.939 & 0.721 & 0.722 & 2.08\\
greenbeb & 0.291 & 0.385 & 0.285 & 1.02\\
cycle & 0.024 & 0.242 & 0.013 & 0.042\\
greenbea & 0.158 & 0.230 & 0.166 & 2.62\\
degen3 & 0.169 & 0.180 & 0.134 & 0.322\\
25fv47 & 0.231 & 0.167 & 0.134 & 0.269\\
pilotnov & 0.095 & 0.158 & 0.078 & 0.235\\
nesm & 0.068 & 0.157 & 0.081 & 0.630\\
fit2d & 0.076 & 0.151 & 0.070 & 0.149\\
pilot.we & 0.130 & 0.139 & 0.130 & 0.374\\
80bau3b & 0.117 & 0.132 & 0.080 & 0.457\\
d6cube & 0.064 & 0.127 & 0.683 & 0.110\\
grow22 & 0.126 & 0.102 & 0.044 & 0.246\\
pilot.ja & 0.129 & 0.102 & 0.178 & 0.836\\
perold & 0.082 & 0.087 & 0.084 & 0.341\\
woodw & 0.031 & 0.078 & 0.097 & 0.157\\
bnl2 & 0.066 & 0.075 & 0.035 & 0.383\\
fit1p & 0.065 & 0.067 & 0.026 & 0.122\\
scsd8 & 0.100 & 0.057 & 0.023 & 0.087\\
stocfor2 & 0.023 & 0.043 & 0.029 & 0.177\\
\bottomrule
\end{tabular}
\end{table}

Table~\ref{tab:kennington} shows the results on the 16 problems of the Kennington test set. All four solvers solved all 16 problems. The geometric mean of ENOMOTO-Solver is almost the same as that of CLP (0.303\,s vs.\ 0.293\,s), and its shifted geometric mean and total time are the smallest of the four. In particular, on the four osa problems it needed 0.09--0.2 times the time of HiGHS and 0.4--0.8 times that of CLP. On the other hand, on ken-18 and pds-20 it is much slower than HiGHS or CLP.

\begin{table}[htbp]
\centering
\caption{Computation times on the Kennington test set (seconds, median of three runs). $m,n$ are the numbers of rows and columns of the original problem, and nnz is the number of nonzeros of the constraint matrix.}
\label{tab:kennington}
\small
\setlength{\tabcolsep}{4pt}
\begin{tabular}{lrrrrrrr}
\toprule
Problem & $m$ & $n$ & nnz & ENOMOTO & HiGHS & CLP & SoPlex\\
\midrule
cre-a & 3516 & 4067 & 14987 & 0.084 & 0.089 & 0.044 & 0.118\\
cre-b & 9648 & 72447 & 256095 & 1.41 & 1.29 & 1.35 & 3.81\\
cre-c & 3068 & 3678 & 13244 & 0.074 & 0.077 & 0.043 & 0.080\\
cre-d & 8926 & 69980 & 242646 & 0.733 & 0.629 & 0.553 & 2.28\\
ken-07 & 2426 & 3602 & 8404 & 0.015 & 0.038 & 0.011 & 0.043\\
ken-11 & 14694 & 21349 & 49058 & 0.156 & 0.238 & 0.130 & 1.27\\
ken-13 & 28632 & 42659 & 97246 & 0.824 & 0.783 & 0.372 & 6.53\\
ken-18 & 105127 & 154699 & 358171 & 7.66 & 4.29 & 2.48 & 128\\
osa-07 & 1118 & 23949 & 143694 & 0.045 & 0.226 & 0.054 & 0.184\\
osa-14 & 2337 & 52460 & 314760 & 0.134 & 0.871 & 0.183 & 1.06\\
osa-30 & 4350 & 100024 & 600138 & 0.476 & 4.12 & 0.675 & 5.50\\
osa-60 & 10280 & 232966 & 1397793 & 1.71 & 18.7 & 4.19 & 36.9\\
pds-02 & 2953 & 7535 & 16390 & 0.031 & 0.100 & 0.023 & 0.066\\
pds-06 & 9881 & 28655 & 62524 & 0.187 & 0.384 & 0.257 & 0.793\\
pds-10 & 16558 & 48763 & 106436 & 0.557 & 0.719 & 0.863 & 2.38\\
pds-20 & 33874 & 105728 & 230200 & 3.38 & 2.33 & 11.2 & 19.6\\
\bottomrule
\end{tabular}
\end{table}

Table~\ref{tab:mittelmann} shows the 22 of the 40 Mittelmann problems that some solver solved within 600\,s; no solver solved the remaining 18. CLP solved the most problems (18), followed by ENOMOTO-Solver (12), HiGHS (11), and SoPlex (7), while ENOMOTO-Solver has the smallest shifted geometric mean. Only ENOMOTO-Solver solved physiciansched3-3 (515\,s), and on neos-5052403-cygnet and fome13 it needed about 0.03 times and at most 0.22 times, respectively, the time of the other solvers that solved them. Conversely, it did not solve neos within the time limit (SoPlex needed 54.7\,s), and on cont1 and square41 it needed 1.6--2.1 times the time of HiGHS.

\begin{table}[htbp]
\centering
\caption{The 22 problems of Mittelmann's LP benchmarks that some solver solved within 600\,s (times in seconds). ``---'' means not solved within the time limit. Values of 300\,s or more are single measurements; the others are medians of three runs.}
\label{tab:mittelmann}
\small
\begin{tabular}{lrrrr}
\toprule
Problem & ENOMOTO & HiGHS & CLP & SoPlex\\
\midrule
cont1 & 476.5 & 226.6 & 483.4 & ---\\
datt256 & --- & --- & 595.3 & ---\\
ex10 & 105.8 & 107.0 & --- & ---\\
fome13 & 16.1 & 77.4 & 72.5 & 355.7\\
irish-electricity & 161.6 & 181.9 & --- & ---\\
Linf\_520c & --- & --- & 478.0 & ---\\
neos & --- & 390.0 & 260.3 & 54.7\\
neos-5052403-cygnet & 16.1 & --- & 554.2 & 521.5\\
neos-5251015 & --- & --- & 470.0 & ---\\
ns1688926 & --- & --- & 18.5 & 112.3\\
nug08-3rd & 342.8 & 274.0 & 485.3 & ---\\
pds-100 & 108.3 & 89.6 & 167.5 & ---\\
physiciansched3-3 & 515.3 & --- & --- & ---\\
qap15 & --- & --- & 475.1 & ---\\
rail4284 & --- & --- & 42.4 & ---\\
rmine15 & --- & --- & 464.3 & ---\\
s250r10 & 82.2 & 104.7 & 363.0 & 502.7\\
scpm1 & --- & --- & 86.0 & ---\\
square41 & 112.5 & 70.2 & --- & 211.9\\
stormG2\_1000 & 79.5 & 70.8 & 192.6 & ---\\
supportcase10 & 207.5 & 208.0 & 105.2 & 163.2\\
woodlands09 & --- & --- & 359.2 & ---\\
\bottomrule
\end{tabular}
\end{table}

\subsection{Infeasible and unbounded problems}
\label{sec:exp-infeasible}

On the 58 infeasible and unbounded problems, ENOMOTO-Solver identified the correct case every time and was faster than HiGHS on 55 of them (geometric mean of the ratios 0.28). As Table~\ref{tab:summary} shows, it is the fastest of the four solvers on the infeasible problems and second to SoPlex on the unbounded ones.

Of the 29 infeasible problems, 17 were declared infeasible during presolve. On the remaining 12, Phase A did not detect $z^1<0$ (on 11 of them no variable was placed on an infinite side in the initial configuration and Phase A was skipped, and on one $z^1=0$), and infeasibility was detected in Phase B.

On 27 of the 29 unbounded problems, ENOMOTO-Solver detected $z^1<0$, that is, the absence of a finite optimal solution, at the end of Phase A, and concluded unboundedness after Phase B and the cleanup (the remaining two are the dual of cplex2, where $z^1=0$, and a problem declared unbounded on a trivial connected component). Table~\ref{tab:phase-times} shows the times on these 27 problems. The detection of $z^1<0$ at the end of Phase A occurred on average at 98\% (at least 88\%) of the time at which unboundedness was concluded, so distinguishing infeasibility from unboundedness took little additional time. This is consistent with the fact that, whereas the existing two-phase dual simplex method needs a separate procedure such as the primal simplex method to distinguish the two cases once the optimal value of phase 1 is negative (Section~\ref{sec:known-two-phase}), our method obtains the distinction as a continuation into Phase B.

\begin{table}[htbp]
\centering
\caption{Times of ENOMOTO-Solver on the 27 unbounded problems on which $z^1<0$ was detected at the end of Phase A (seconds from the start of the solve).}
\label{tab:phase-times}
\small
\begin{tabular}{lrr}
\toprule
Point & Geo.\ mean & Total\\
\midrule
End of Phase A (detection of $z^1<0$) & 0.00786 & 1.719\\
Conclusion of unboundedness (end of Phase B and cleanup) & 0.00801 & 1.722\\
Whole solve (including postsolve) & 0.00807 & 1.724\\
\bottomrule
\end{tabular}
\end{table}

As shown above, ENOMOTO-Solver is faster than HiGHS at the scale of Netlib and Kennington, and its shifted geometric mean and total time are the smallest of the four solvers. On the large Mittelmann problems, it solves fewer problems than CLP, and on some problems, such as cont1, neos, square41, ken-18, and pds-20, it remains much slower than the other solvers. For large sparse problems, the right-hand sides and solution vectors of FTRAN and BTRAN, as well as the pivot row, become extremely sparse (hyper-sparse), and operations that exploit this sparsity largely determine the computation time\cite{hall2005}. Analyzing the causes of the differences on the slow problems is left for future work, but these linear algebra techniques are independent of the method and can be applied to each phase as they are.

\section*{Statements and declarations}

\paragraph{Funding.} No funding was received for conducting this study.

\paragraph{Competing interests.} The author has no competing interests to declare that are relevant to the content of this article.

\paragraph{Data availability.} The Netlib, Kennington and Mittelmann test problems are publicly available from their original sources \cite{gay1985,carolan1990,mittelmann}. The unbounded instances constructed as duals of the infeasible problems, the measurement scripts and the raw results of all experiments are available at \url{https://github.com/enomotokan/enomoto_solver}.

\paragraph{Code availability.} The source code of ENOMOTO-Solver is available at \url{https://github.com/enomotokan/enomoto_solver} under the MIT License.

\paragraph{Use of generative AI.} In this work, we used Claude, a generative AI by Anthropic (used via Claude Code), for the following purposes.
\begin{itemize}
\item \textbf{Solver implementation}: most of the implementation of ENOMOTO-Solver was carried out by the AI under the author's design and direction. Many of the implementation techniques listed in Section~\ref{sec:exp-features} were proposed by the AI, and the author decided whether to adopt them based on the results of comparative experiments.
\item \textbf{Numerical experiments}: the AI was used to write the measurement scripts, run the measurements, and aggregate the results.
\item \textbf{Manuscript preparation}: the AI was used for drafting the text of the manuscript, organizing the written proofs, surveying the literature, and revising.
\end{itemize}
The ideas and theory of the method in this paper (the analysis with a symbolic $M$, the decomposition into the slope problem and the intercept problem, the distinction of the three cases, and the correspondence between the big-$M$ method and the two-phase method for the dual problem) are due to the author. All text, proofs, code, and numerical results generated by the AI have been checked by the author, and the author takes responsibility for the content of this paper.

\begin{thebibliography}{99}

\bibitem{dantzig1963}
G.~B.~Dantzig,
\textit{Linear Programming and Extensions},
Princeton University Press, 1963.

\bibitem{dantzig1955}
G.~B.~Dantzig, A.~Orden, P.~Wolfe,
``The generalized simplex method for minimizing a linear form under linear inequality restraints,''
\textit{Pacific Journal of Mathematics}, 5(2), 183--195, 1955.

\bibitem{lemke1954}
C.~E.~Lemke,
``The dual method of solving the linear programming problem,''
\textit{Naval Research Logistics Quarterly}, 1(1), 36--47, 1954.

\bibitem{bland1977}
R.~G.~Bland,
``New finite pivoting rules for the simplex method,''
\textit{Mathematics of Operations Research}, 2(2), 103--107, 1977.

\bibitem{forrestgoldfarb1992}
J.~J.~Forrest, D.~Goldfarb,
``Steepest-edge simplex algorithms for linear programming,''
\textit{Mathematical Programming}, 57(1--3), 341--374, 1992.

\bibitem{harris1973}
P.~M.~J.~Harris,
``Pivot selection methods of the Devex LP code,''
\textit{Mathematical Programming}, 5(1), 1--28, 1973.

\bibitem{maros2003}
I.~Maros,
\textit{Computational Techniques of the Simplex Method},
Springer, 2003.

\bibitem{chvatal1983}
V.~Chv\'atal,
\textit{Linear Programming},
W.~H.~Freeman, 1983.

\bibitem{bertsimas1997}
D.~Bertsimas, J.~N.~Tsitsiklis,
\textit{Introduction to Linear Optimization},
Athena Scientific, 1997.

\bibitem{sherali1983}
H.~D.~Sherali, A.~L.~Soyster,
``Preemptive and nonpreemptive multi-objective programming: Relationship and counterexamples,''
\textit{Journal of Optimization Theory and Applications}, 39(2), 173--186, 1983.

\bibitem{cococcioni2021}
M.~Cococcioni, L.~Fiaschi,
``The Big-M method with the numerical infinite M,''
\textit{Optimization Letters}, 15(7), 2455--2468, 2021.

\bibitem{wolfe1965}
P.~Wolfe,
``The composite simplex algorithm,''
\textit{SIAM Review}, 7(1), 42--54, 1965.

\bibitem{fourer1994}
R.~Fourer,
``Notes on the dual simplex method,''
Draft report, Northwestern University, 1994.

\bibitem{kostina2002}
E.~Kostina,
``The long step rule in the bounded-variable dual simplex method: Numerical experiments,''
\textit{Mathematical Methods of Operations Research}, 55(3), 413--429, 2002.

\bibitem{maros2003pl}
I.~Maros,
``A piecewise linear dual phase-1 algorithm for the simplex method,''
\textit{Computational Optimization and Applications}, 26(1), 63--81, 2003.

\bibitem{koberstein2005}
A.~Koberstein,
\textit{The Dual Simplex Method, Techniques for a Fast and Stable Implementation},
Dissertation, Universit\"at Paderborn, 2005.

\bibitem{bixby2002}
R.~E.~Bixby,
``Solving real-world linear programs: A decade and more of progress,''
\textit{Operations Research}, 50(1), 3--15, 2002.

\bibitem{koberstein2007}
A.~Koberstein, U.~H.~Suhl,
``Progress in the dual simplex method for large scale LP problems: practical dual phase 1 algorithms,''
\textit{Computational Optimization and Applications}, 37(1), 49--65, 2007.

\bibitem{huang2021}
M.~Huang, Y.~Zhong, H.~Yang, J.~Wang, F.~Zhang, B.~Bai, L.~Shi,
``Simplex initialization: A survey of techniques and trends,''
arXiv:2111.03376, 2021.

\bibitem{megiddo1991}
N.~Megiddo,
``On finding primal- and dual-optimal bases,''
\textit{ORSA Journal on Computing}, 3(1), 63--65, 1991.

\bibitem{bixby1994}
R.~E.~Bixby, M.~J.~Saltzman,
``Recovering an optimal LP basis from an interior point solution,''
\textit{Operations Research Letters}, 15, 169--178, 1994.

\bibitem{zionts1969}
S.~Zionts,
``The criss-cross method for solving linear programming problems,''
\textit{Management Science}, 15(7), 426--445, 1969.

\bibitem{terlaky1985}
T.~Terlaky,
``A convergent criss-cross method,''
\textit{Optimization}, 16(5), 683--690, 1985.

\bibitem{fukuda1997}
K.~Fukuda, T.~Terlaky,
``Criss-cross methods: A fresh view on pivot algorithms,''
\textit{Mathematical Programming}, 79(1--3), 369--395, 1997.

\bibitem{vanderbei2020}
R.~J.~Vanderbei,
\textit{Linear Programming: Foundations and Extensions}, 5th ed.,
International Series in Operations Research \& Management Science, vol.~285,
Springer, 2020.

\bibitem{huangfu2018}
Q.~Huangfu, J.~A.~J.~Hall,
``Parallelizing the dual revised simplex method,''
\textit{Mathematical Programming Computation}, 10(1), 119--142, 2018.

\bibitem{hall2005}
J.~A.~J.~Hall, K.~I.~M.~McKinnon,
``Hyper-sparsity in the revised simplex method and how to exploit it,''
\textit{Computational Optimization and Applications}, 32(3), 259--283, 2005.

\bibitem{gilbert1988}
J.~R.~Gilbert, T.~Peierls,
``Sparse partial pivoting in time proportional to arithmetic operations,''
\textit{SIAM Journal on Scientific and Statistical Computing}, 9(5), 862--874, 1988.

\bibitem{forrest1972}
J.~J.~H.~Forrest, J.~A.~Tomlin,
``Updated triangular factors of the basis to maintain sparsity in the product form simplex method,''
\textit{Mathematical Programming}, 2(1), 263--278, 1972.

\bibitem{markowitz1957}
H.~M.~Markowitz,
``The elimination form of the inverse and its application to linear programming,''
\textit{Management Science}, 3(3), 255--269, 1957.

\bibitem{ruiz2001}
D.~Ruiz,
``A scaling algorithm to equilibrate both rows and columns norms in matrices,''
Technical Report RAL-TR-2001-034, Rutherford Appleton Laboratory, 2001.

\bibitem{gay1985}
D.~M.~Gay,
``Electronic mail distribution of linear programming test problems,''
\textit{Mathematical Programming Society COAL Newsletter}, 13, 10--12, 1985.

\bibitem{gurobi}
Gurobi Optimization, LLC,
\textit{Gurobi Optimizer Reference Manual},
\url{https://docs.gurobi.com/projects/optimizer/en/current/} (accessed September 25, 2026).
See the status codes INF\_OR\_UNBD and UNBOUNDED, the parameters DualReductions and Method, and the section on Concurrent Optimization.

\bibitem{cplex}
IBM,
\textit{IBM ILOG CPLEX Optimization Studio 20.1.0 Documentation},
\url{https://www.ibm.com/docs/en/icos/20.1.0} (accessed September 25, 2026).
See the solution status codes (CPX\_STAT\_INForUNBD) and the parameter LPMETHOD.

\bibitem{carolan1990}
W.~J.~Carolan, J.~E.~Hill, J.~L.~Kennington, S.~Niemi, S.~J.~Wichmann,
``An empirical evaluation of the KORBX algorithms for military airlift applications,''
\textit{Operations Research}, 38(2), 240--248, 1990.

\bibitem{mittelmann}
H.~D.~Mittelmann,
``Benchmarks for optimization software,''
\url{https://plato.asu.edu/bench.html}.

\bibitem{clp}
J.~J.~Forrest, D.~de~la~Nuez, R.~Lougee-Heimer,
\textit{CLP User Guide},
COIN-OR, 2004.
\url{https://github.com/coin-or/Clp} (version 1.17.11 was used).

\bibitem{soplex}
R.~Wunderling,
\textit{Paralleler und objektorientierter Simplex-Algorithmus},
Dissertation, Technische Universit\"at Berlin, 1996.
\url{https://github.com/scipopt/soplex} (version 8.1.0 was used).

\bibitem{andersen1996}
E.~D.~Andersen, Y.~Ye,
``Combining interior-point and pivoting algorithms for linear programming,''
\textit{Management Science}, 42(12), 1719--1731, 1996.

\end{thebibliography}

\end{document}
